\BourbakiExerciseGroup

\begin{enumerate}
\def\labelenumi{\arabic{enumi})}
\tightlist
\item
  设 \(U = (\alpha_{ij})\) 是主理想整环 A 上一个 \(m\) 行 \(n\) 列的矩阵。a) 首先假设诸 \(\alpha_{ij}\) 整体互素。证明存在两个元素属于 A 的可逆方阵 \(P\) 和 \(Q\)，使得矩阵 \(PUQ\) 的某一元素等于 1。（对每个 \(\alpha_{ij} \neq 0\) ，设 \(S(\alpha_{ij})\) 为 \(\alpha_{ij}\) 分解为不可约元时诸指数之和，并设 \(s(U)\) 为诸数 \(s(\alpha_{ij})\) 中最小者；若 \(s(U) > 0\) ，证明存在两个可逆方阵 \(R\) 和 \(S\)，使得
\end{enumerate}

\[
s (R U S) <   s (U);
\]

利用贝祖恒等式，并在必要时对行和列作置换后，把所考虑的矩阵 \(R\) 和 \(S\) 限于形如 \(\begin{pmatrix} T & 0 \\ 0 & I \end{pmatrix}\) 的矩阵，其中 \(T\) 是一个 \(2 \times 2\) 矩阵，\(I\) 是单位矩阵，或限于此类矩阵与置换矩阵的乘积。）

\begin{enumerate}
\def\labelenumi{\alph{enumi})}
\setcounter{enumi}{1}
\tightlist
\item
  若 \(\delta_1\) 是 \(U\) 的诸元素的一个最大公因子，证明存在两个可逆矩阵 \(P_1, Q_1\)，使得
\end{enumerate}

\[
P _ {1} U Q _ {1} = \left( \begin{array}{c c} \delta_ {1} & 0 \\ 0 & U _ {1} \end{array} \right)
\]

其中 \(U_{1}\) 的所有元素都可被 \(\delta_{1}\) 整除（利用 a））。

\begin{enumerate}
\def\labelenumi{\alph{enumi})}
\setcounter{enumi}{2}
\tightlist
\item
  对于 \(\mathbf{A} = \mathbf{Z}\) ，利用 \(b\) 得到一个计算元素属于 \(\mathbf{Z}\) 的具体矩阵的不变因子的方法。将此方法应用于矩阵
\end{enumerate}

\[
\left( \begin{array}{c c c c} 6 & 8 & 4 & 20 \\ 12 & 12 & 18 & 30 \\ 18 & 4 & 4 & 10 \end{array} \right).
\]

\begin{enumerate}
\def\labelenumi{\arabic{enumi})}
\setcounter{enumi}{1}
\item
  设 A 为主理想整环。则 \(A^{q}\) 的两个子模 M 与 N 在 \(A^{q}\) 的自同构下互为变换，当且仅当它们关于 \(A^{q}\) 具有相同的不变因子。
\item
  设 A 为主理想整环，其分式域为 K，设 E 为 K 上的向量空间，设 M 为包含于 E 中的秩为 n 的有限生成 A-模，并设 N 为包含于 KM 中的秩为 \(p \leqslant n\) 的有限生成 A-模。证明存在 M 的一组基 \((e_{i})_{1 \leqslant i \leqslant n}\) 以及 N 的一组由向量 \(\alpha_{i}e_{i}\) ( \(1 \leqslant i \leqslant p\) ) 组成的基，其中 \(\alpha_{i} \in K\) ，且 \(\alpha_{i}\) （关于环 A）整除 \(\alpha_{i+1}\) ，对 \(1 \leqslant i \leqslant p - 1\) ；同时证明分式理想 \(A\alpha_{i}\) 是唯一确定的（即 N 关于 M 的不变因子）。
\item
  设 \(G\) 是一个有限阿贝尔群。
\end{enumerate}

\begin{enumerate}
\def\labelenumi{\alph{enumi})}
\item
  对每个素数 \(p\) ，证明阶为 \(p\) 的元素在 \(G\) 中的个数是 \(p^N - 1\) ，其中 \(N = \sum_{n=1}^{\infty} m(p^n)\) （采用 VII，第 24 页，命题 9 的记号）。
\item
  由此推出，若方程 \(x^p = 1\) 对每个素数 \(p\) 在 \(G\) 中至多有 \(p\) 个解，则 \(G\) 是循环群。
\end{enumerate}

\begin{enumerate}
\def\labelenumi{\arabic{enumi})}
\setcounter{enumi}{4}
\item
  设 \(G\) 为一个 \(n\) 阶有限交换群。证明：只要整数 \(q\) 整除 \(n\) ，就存在一个 \(q\) 阶子群包含在 \(G\) 中（把 \(G\) 分解为不可分解循环子群的直和）。
\item
  设 A 是一个交换环，并设 E 是一个 A-模。设 \(\mathcal{E}\) 表示（无算子的）阿贝尔群 E 的自同态环 \(\operatorname{End}_{\mathbf{z}}(\mathbf{E})\) ；对每个子环 \(\mathbf{B} \subset \mathcal{E}\) ，设 \(\mathbf{B}'\) 表示 B 在 \(\mathcal{E}\) 中的中心化子（即 \(\mathcal{E}\) 中由与 B 的每个元素都可交换的元素组成的子环）。
\end{enumerate}

\begin{enumerate}
\def\labelenumi{\alph{enumi})}
\item
  设 \(A_{0}\) 为 E 的由位似变换 \(x \mapsto \lambda x (\lambda \in A)\) 组成的子环；若 a 是 E 的零化理想，则已知 \(A_{0}\) 同构于 A/a，且它在 E 中的中心化子 \(A_{0}^{\prime}\) 是 E 作为 A-模的自同态环 \(\operatorname{End}_{\mathbf{A}}(\mathbf{E})\) 。证明中心化子 \(A_{0}^{\prime\prime}\) （\(A_{0}^{\prime}\) 的中心化子）是交换的，且其中心化子 \(A_{0}^{\prime\prime\prime}\) 等于 \(A_{0}^{\prime}\) （注意 \(A_{0} \subset A_{0}^{\prime}\) ）。
\item
  如果子模 F 是 E 的一个直因子，证明 \(u(F) \subset F\) 对所有 \(u \in A_{0}''\) （考虑由 E 分解为 F 与另一个子模的直和所得到的从 E 到 F 的投影）。
\item
  假设 E 是循环子模的一个族 \((F_{i})\) 的直和。若 \(u \in A_{0}''\) ，证明对每个指标 i，存在元素 \(\alpha_{i} \in A\) ，使得 \(u(x) = \alpha_{i}x\) 对所有 \(x \in F_{i}\) 成立（利用 b））。
\item
  假设 E 是循环子模的一个序列 \((\mathrm{E}_{n})\) （有限或无限）的直和，使得：若 \(b_{n}\) 表示 \(E_{n}\) 的零化理想，则序列 \((\mathfrak{b}_{n})\) 是递增的。证明 \(A_{0}^{\prime\prime}=A_{0}\) （利用 c)，以及对每个指标 i\textgreater1 存在一个自同态 \(v_{i}\in A_{0}^{\prime}\) 把 \(E_{i-1}\) 映入 \(E_{i}\) 这一事实。）
\end{enumerate}

将此应用于 A 为主理想整环且 E 为有限生成 A-模的情形。

\begin{enumerate}
\def\labelenumi{\arabic{enumi})}
\setcounter{enumi}{6}
\tightlist
\item
  设 A 为主理想整环，其分式域为 K，并设 E 为秩 n 的有限生成无挠 A-模。设 \(E^{*}\) 为 E 的对偶；则 E 与 \(E^{*}\) 均同构于 \(A^{n}\) .
\end{enumerate}

\begin{enumerate}
\def\labelenumi{\alph{enumi})}
\item
  设 M 为 E 的子模；证明：若 \(M^{0}\) 是 \(E^{*}\) 中与 M 正交的子模，则模 \(E^{*}/M^{0}\) 是无挠的，从而 \(M^{0}\) 是 \(E^{*}\) 的直因子；子模 \(M^{00}\) （E 中与 \(M^{0}\) 正交者）等于 \(E \cap KM\) （把 E 视为典范地嵌入到 K 上一个 n 维向量空间中）。可把 \(E^{*}/M^{0}\) 视为典范地嵌入到 M 的对偶 \(M^{*}\) 中；证明 \(E^{*}/M^{0}\) 关于 \(M^{*}\) 的不变因子等于 M 关于 E 的不变因子（利用 VII 第 18 页定理 1）。
\item
  设 F 是第二个有限生成的无挠 A-模，并设 u 是从 E 到 F 的线性变换。证明 \(u(F^{*})\) 关于 \(E^{*}\) 的不变因子与 \(u(E)\) 关于 F 的那些不变因子相等（按 VII 第 21 页命题 5 的方式进行）。
\end{enumerate}

\begin{enumerate}
\def\labelenumi{\arabic{enumi})}
\setcounter{enumi}{7}
\tightlist
\item
  设 G 是主理想整环 A 上的一个有限长度模。对于满足 \(N \subset M\) 且使 G 同构于 M/N 的任意一对自由模 M 和 N，N 关于 M 的不变因子中异于 A 的那些与 G 的不变因子相同（VII，第 20 页，定义 2）；它们的个数称为 G 的秩。
\end{enumerate}

\begin{enumerate}
\def\labelenumi{\alph{enumi})}
\item
  证明 G 的秩 r 是使 G 为其直和的循环子模的最少个数，且等于 G 的诸 \(\pi\) -准素分量的秩中的最大者。b) 证明 G 的子模或商模的秩至多等于 G 的秩（把 G 视为两个秩为 r 的自由模的商）。
\item
  对每个 \(\lambda \in \mathbf{A}\) ，证明子模 \(\lambda G\) 的秩等于 \(G\) 的不变因子中不整除 \(\lambda\) 的那些的个数（注意，若 \(E = A / (A\alpha)\) ，则模 \(\lambda E\) 同构于 \((A\lambda) / ((A\alpha) \cap (A\lambda))\) ）。由此推出，第 \(k\) 个不变因子（按递减顺序，属于 \(G\) ）是那些 \(\lambda \in A\) 的一个最大公因子，其中 \(\lambda G\) 的秩 \(\leqslant r - k\) 。
\item
  设 \(A\alpha_{k}\) ( \(1 \leqslant k \leqslant r\) ) 是 G 的不变因子，按递减顺序排列。设 H 是 G 的秩为 r - q 的子模，并设 \(A\beta_{k}\) ( \(1 \leqslant k \leqslant r - q\) ) 是它的不变因子，按递减顺序排列。由 b) 和 c) 推出 \(\beta_{k}\) 整除 \(\alpha_{k+q}\) ，对 \(1 \leqslant k \leqslant r - q\) 。类似地证明，若 G/H 的秩为 r - p 且不变因子为 \(A\gamma_{k}\) ( \(1 \leqslant k \leqslant r - p\) )（按递减顺序），则 \(\gamma_{k}\) 整除 \(\alpha_{k+p}\) ，对 \(1 \leqslant k \leqslant r - p\) 。
\item
  反之，设 L 是秩为 r - q 的挠模，其不变因子 \(A\lambda_{k}\) ( \(1 \leqslant k \leqslant r - q\) ) 满足 \(\lambda_{k}\) 整除 \(\alpha_{k+q}\) ，对 \(1 \leqslant k \leqslant r - q\) 。证明存在 G 的子模 M 和 N，使 L 同构于 M 且同构于 G/N（把 G 分解为同构于 A/A \(\alpha_{k}\) 的子模的直和）。
\end{enumerate}

\begin{enumerate}
\def\labelenumi{\arabic{enumi})}
\setcounter{enumi}{8}
\tightlist
\item
  设 A 是一个主理想整环，K 为其分式域；设 E 是 K 上的向量空间，设 M 是包含在 E 中的秩为 n 的有限生成 A-模，设 N 是包含在 KM 中的秩为 p 的有限生成 A-模，设 P 是包含在 KN 中的秩为 q 的有限生成 A-模。
\end{enumerate}

\begin{enumerate}
\def\labelenumi{\alph{enumi})}
\item
  设 \(\mathbf{A}\alpha_{i}\) ( \(1 \leqslant i \leqslant p\) ) 是 \(\mathbf{N}\) 关于 \(\mathbf{M}\) 的不变因子，按递减顺序排列（习题 3）。证明，对 \(1 \leqslant k \leqslant p\) ，元素 \(\alpha_{k}\) 是那些元素 \(\lambda \in \mathbf{K}\) 的一个最大公因子，这些元素使商模 \((\mathbf{N} + \lambda (\mathbf{KN} \cap \mathbf{M})) / \mathbf{N}\) 的秩 \(\leqslant p - k\)（参见习题 8，d））。
\item
  设 \(\mathbf{A}\beta_{j}\) ( \(1 \leqslant j \leqslant q\) ) 是 \(\mathbf{P}\) 关于 \(\mathbf{M}\) 的不变因子，\(\mathbf{A}\gamma_{j}\) ( \(1 \leqslant j \leqslant q\) ) 是 \(\mathbf{P}\) 关于 \(\mathbf{N}\) 的不变因子，均按递减顺序排列。证明 \(\alpha_{1}\gamma_{j}\) 整除 \(\beta_{j}\) ，对 \(1 \leqslant j \leqslant q\) （参见习题 8）。
\item
  证明 \(\gamma_{1}\alpha_{j}\) 整除 \(\beta_{j}\) ，对 \(1 \leqslant j \leqslant q\) 。（首先考虑 \(n = p = q\) 的情形，并应用习题 8。一般地，证明我们总可以假设 \(p = n\) ，并考虑 M 的一个子模 Q，使得 \(\mathbf{P} + \rho \mathbf{Q}\) 的秩为 \(p\) ，对所有 \(\rho \in \mathbf{A}\) ；然后选取足够小的理想 \(\mathbf{A}\rho\) ，并应用 VII 第 20 页命题 4。）
\end{enumerate}

\(d\) ) 对 N 的任意秩为 \(k \leqslant p\) 的子模 H，设 \((\mu_{\mathrm{H}})\) 为由那些 \(\mu \in \mathbf{K}\) 组成的分式理想，它们满足 \(\mu (\mathbf{M} \cap \mathbf{KH}) \subset \mathbf{H}\) 。由 \(c\) ) 推出，对每个 \(1 \leqslant k \leqslant p\) ，元素 \(\alpha_{k}\) 是诸 \(\mu_{\mathrm{H}}\) 的一个最大公因子，这里 H 遍历 N 的所有秩为 \(k\) 的子模。

\begin{enumerate}
\def\labelenumi{\arabic{enumi})}
\setcounter{enumi}{9}
\tightlist
\item
  设 A 为主理想整环，M 是 \(E = A^{n}\) 中秩为 n 的子模；设 \(A\alpha_{i} (1 \leqslant i \leqslant n)\) 为 M 关于 E 的不变因子，按递减顺序排列。设 N 是 M 的在 M 中有补的子模（即满足 \(N = M \cap KN\) ，其中 K 是 A 的分式域）。
\end{enumerate}

\begin{enumerate}
\def\labelenumi{\alph{enumi})}
\item
  设 P 是 N 在 M 中的一个补。证明 E/M 的子模 \((P + (E \cap KN))/M\) 同构于 \((E \cap KN)/N\) 。由此推出，若 N 的秩为 p 且 \(A\beta_{k}\) ( \(1 \leqslant k \leqslant p\) ) 是 N 关于 E 的不变因子，则 \(\beta_{k}\) 整除 \(\alpha_{k+n-p}\) ，对 \(1 \leqslant k \leqslant p\) （利用习题 8，d））。
\item
  给出一个例子，说明模 E/M 不必同构于模 \((E \cap KN)/N\) 与 \(E/(P + (E \cap KN))\) 的乘积（取 E/M 为不可分解的）。
\end{enumerate}

\begin{enumerate}
\def\labelenumi{\arabic{enumi})}
\setcounter{enumi}{10}
\item
  设 A 为主理想整环，H 是 \(E = A^{n}\) 的一个子模，\(H^{0}\) 是 E 的对偶 \(E^{*}\) 中与 H 正交的子模（习题 7）。设 M 是 E 的一个秩为 p 的子模，并设 \(A\alpha_{i}\) ( \(1 \leqslant i \leqslant p\) ) 为 M 关于 E 的不变因子，按递减顺序排列。设 \((\nu_{\mathrm{H}})\) 为 A 的诸理想 \(f(\mathbf{M})\) 的最大公因子，这里 f 遍历 \(H^{0}\) ；证明，对 \(1 \leqslant k \leqslant p\) ，元素 \(\alpha_{k}\) 是 E 的秩为 k - 1 且在 E 中有补的所有子模 H 所对应的 \(\nu_{H}\) 的一个最大公因子（首先考虑 k = 1 的情形；把结果应用于 E/H 的子模 \((\mathbf{M} + \mathbf{H})/\mathbf{H}\) ，并对 E/M 的商模 \(E/(M + H)\) 利用习题 8，d））。
\item
  设 G 是环 A 上的有限长度模。证明：G 的子模 H 若与 G 同构，则必有 H = G（参见 II，第 212 页，命题 16）。
\item
  设 A 为主理想整环，其分式域为 K。
\end{enumerate}

\begin{enumerate}
\def\labelenumi{\alph{enumi})}
\item
  证明：对于每个A-模M，典范同态 \(c_{\mathbf{M}}: \mathbf{M} \to \mathbf{D}(\mathbf{D}(\mathbf{M}))\) 是单射的（参见II，第389页，习题13）。
\item
  若 \(u: \mathbf{M} \to \mathbf{N}\) 是 A-模的一个同态，定义从 \(D(\operatorname{Im}(u))\) 到 \(\operatorname{Im}(D(u))\) 、从 \(D(\operatorname{Ker}(u))\) 到 \(\operatorname{Coker}(D(u))\) 以及从 \(D(\operatorname{Coker}(u))\) 到 \(\operatorname{Ker}(D(u))\) 的自然同构。
\item
  设 M 是一个 A-模，N 是 M 的一个子模。在第 9 节的记号下，证明 \(N^{00} = N\) （利用 a) 和 b)）。模 \(N^{0}\) 具有有限长度，当且仅当 M/N 具有有限长度。反之，对 D(M) 的任意有限长度子模 \(N'\) ，模 \(N'^{0}\) 是 M 的一个子模，使 \(M/N^{0}\) 具有有限长度，且 \(N' = N'^{00}\) 。
\end{enumerate}

\(d\) ) 若 \(\mathbf{M}\) 是挠 A-模，且 \(\mathbf{M}_{\pi}\) 表示 \(\pi\) -准素分量（属于 \(\mathbf{M}\) ），证明 \(\mathrm{D}(\mathbf{M})\) 同构于 \(\prod_{\pi \in \mathbb{P}}\mathrm{Hom}(\mathbf{M}_{\pi},\mathbf{U}_{\pi})\) ，其中 \(\mathbf{P}\) 是 A 的不可约元素的一个代表系，且 \(U_{\pi}\) 是 K/A 的 \(\pi\) -准素分量，对每个 \(\pi \in P\) （VII，第 54 页，习题 3）。

\begin{enumerate}
\def\labelenumi{\alph{enumi})}
\setcounter{enumi}{4}
\tightlist
\item
  由d)推出， \(\mathrm{D}(\mathrm{K}/\mathrm{A})\) 同构于乘积 \(\prod_{\pi\in P}\hat{\mathbf{A}}_{(\pi)}\) ，其中
\end{enumerate}

\(\hat{\mathbf{A}}_{(\pi)}\) 表示 VII，第 54 页，习题 2 中定义的环 \(\mathbf{A}_{(\pi)}\) 的完备化，其拓扑由取理想 \(\pi^n\mathbf{A}_{(\pi)}\) 为 0 的开邻域基来定义（此完备化还可等同于逆向极限 \(\lim (\mathbf{A} / \mathbf{A}\pi^n)\) ）。

\begin{enumerate}
\def\labelenumi{\alph{enumi})}
\setcounter{enumi}{5}
\item
  若 M 是有限长度的 A-模，且 a 是其零化子，证明 D(M) 同构于与 M 相伴的忠实 (A/a)-模的对偶。
\item
  设 M 是有限长度的 A-模。证明：若 N 是 M 的子模，则存在 M 的子模 P，使得 M/N 同构于 P，且 M/P 同构于 N（利用 VII，第 26 页，命题 10 及 VII，第 27 页，命题 12）；具有这些性质的 M 的两个子模称为互反的。证明：对所有 \(\alpha \in A\) ，M 的子模 \(M(\alpha)\) （由所有满足 \(x \in M\) 且 \(\alpha x = 0\) 的元素组成）与子模 \(\alpha M\) 是互反的。
\end{enumerate}

¶ 14) 设 M 是主理想整环 A 上的有限长度模，A 的分式域为 K。

\begin{enumerate}
\def\labelenumi{\alph{enumi})}
\tightlist
\item
  证明：若 N 是 M 的子模，则 M 的秩至多等于 N 与 M/N 的秩之和（令 M = E/H，其中 E 是秩为 n 的自由模，H 是 E 的秩为 n 的子模；若 N = L/H，其中 H ⊆ L ⊆ E，注意，若 N 的秩为 p，则存在 H 的秩为 n - p 的子模 R，使得 R 在 L 中有补 S，且 S ∩ H 是 R 在 H 中的补；此外，(E ∩ KR)/R 的秩 ≥ n - p，且同构于 E/L 的一个商模（习题 10，a））；利用习题 8，b）完成证明。 b) 设 N 是 M 的秩为 p 的子模，q 是 M/N 的秩。设 Aα \(_{i}\) （1 ≤ i ≤ n）为 M 的不变因子，Aβ \(_{k}\) （1 ≤ k ≤ p）为 N 的不变因子，均按递减顺序排列。证明 β \(_{k}\) 是 α \(_{k- (p + q - n)}\) （利用a）及习题8，c），注意到(λM)/(λN)同构于M/N的一个商模）。
\end{enumerate}

\begin{enumerate}
\def\labelenumi{\arabic{enumi})}
\setcounter{enumi}{14}
\tightlist
\item
  \begin{enumerate}
  \def\labelenumii{\alph{enumii})}
  \tightlist
  \item
    设 M 是主理想整环 A 上秩为 n 的有限长度模，并设 \(\mathbf{A}\alpha_{i}\) ( \(1 \leqslant i \leqslant n\) ) 为其按递减顺序排列的不变因子。设 \((\beta_{i})\) ( \(1 \leqslant i \leqslant n\) ) 是 A 中元素的一个序列，使得 \(\beta_{i}\) 整除 \(\alpha_{i}\) ，对 \(1 \leqslant i \leqslant n\) ，且 \(\alpha_{i}\) 整除 \(\beta_{i+1}\) ，对 \(1 \leqslant i \leqslant n-1\) 。设 \((\mathbf{M}_{i})_{1 \leqslant i \leqslant n}\) 是 M 的子模的一个序列，使得 M 是诸 \(M_{i}\) 的直和，且每个 \(M_{i}\) 同构于 \(A/A\alpha_{i}\) ；设 \(a_{i}\) 是 \(M_{i}\) 的一个生成元，对 \(1 \leqslant i \leqslant n\) 。令
  \end{enumerate}
\end{enumerate}

\[
b = \beta_ {1} a _ {1} + \frac {\beta_ {1} \beta_ {2}}{\alpha_ {1}} a _ {2} + \frac {\beta_ {1} \beta_ {2} \beta_ {3}}{\alpha_ {1} \alpha_ {2}} a _ {3} + \dots + \frac {\beta_ {1} \beta_ {2} \dots \beta_ {n}}{\alpha_ {1} \alpha_ {2} \dots \alpha_ {n - 1}} a _ {n}.
\]

证明 M 关于循环子模 Ab 的商模同构于 n 个模 \(A/A\beta_{i}\) ( \(1 \leq i \leq n\) ) 的直和。

\begin{enumerate}
\def\labelenumi{\alph{enumi})}
\setcounter{enumi}{1}
\tightlist
\item
  设 M 和 N 是 A 上两个有限长度模；设 \(A\alpha_{i}\) ( \(1 \leqslant i \leqslant n\) ) 和 \(A\beta_{k}\) ( \(1 \leqslant k \leqslant p\) ) 分别是 M 和 N 的不变因子，按递减顺序排列；假设 \(p \leqslant n\) ，\(\beta_{k}\) 整除 \(\alpha_{k+n-p}\) ，对 \(1 \leqslant k \leqslant p\) ，且 \(\alpha_{k}\) 整除 \(\beta_{k+(p+q-n)}\) ，对 \(1 \leqslant k \leqslant n-q\) ，其中 q 是一个整数，使得
\end{enumerate}

\[
q \leqslant n \leqslant p + q.
\]

证明存在 M 的一个与 N 同构的子模 P，使得 M/P 的秩 \(\leq q\) （把两个模 M 和 N 各自分解为 q 个模的直和，从而把问题归结为 q = 1 的情形；然后利用 a) 以及习题 13 的 g)）。

\begin{enumerate}
\def\labelenumi{\alph{enumi})}
\setcounter{enumi}{2}
\tightlist
\item
  对于 A 上的模 N，要使它与有限长度模 M 的一个直和因子同构，当且仅当 N 的每个初等因子都是 M 的初等因子，且其在 N 中的重数至多等于其在 M 中的重数。由此给出一个有限长度模 M 及其子模 N 的例子，使得 M 的秩等于 N 与 M/N 的秩之和，但 N 在 M 中没有补（利用 a)）。
\end{enumerate}

\begin{enumerate}
\def\labelenumi{\arabic{enumi})}
\setcounter{enumi}{15}
\tightlist
\item
  设 M 是交换环 A 上的一个模。若对 M 的每个自同态 u 都有 \(u(N) \subset N\) ，则 M 的子模 N 称为特征子模。对所有 \(\alpha \in A\) ，子模 \(M(\alpha)\) （由被 \(\alpha\) 零化的元素组成）与子模 \(\alpha M\) 都是特征的。
\end{enumerate}

\begin{enumerate}
\def\labelenumi{\alph{enumi})}
\item
  若 N 是 M 的特征子模，且 P 与 Q 是 M 的两个互补子模，证明 N 是 \(N \cap P\) 与 \(N \cap Q\) 的直和（考虑从 M 到 P 和 Q 上的投影）。
\item
  设 M 是主理想整环 A 上的一个有限长度模；设 \(A\alpha_{i}\) 为 M 的不变因子，按递减顺序排列 \((1 \leqslant i \leqslant n)\) ；设 \((\mathbf{M}_{i})_{1 \leqslant i \leqslant n}\) 是 M 的子模的一个序列，使得 M 是诸 \(M_{i}\) 的直和，且每个 \(M_{i}\) 同构于 \(A/A\alpha_{i}\) 。设 N 是 M 的一个特征子模，并设 \(A\beta_{i}\) 为 \(N \cap M_{i}\) 的零化理想 \((1 \leqslant i \leqslant n)\) ；若 \(\alpha_{i} = \beta_{i}\gamma_{i}\) ，证明 \(\beta_{i}\) 整除 \(\beta_{i+1}\) 且 \(\gamma_{i}\) 整除 \(\gamma_{i+1}\) ，对所有 \(1 \leqslant i \leqslant n-1\) （注意，对 \(1 \leqslant i \leqslant n-1\) ，存在 M 的一个把 \(M_{i+1}\) 映到 \(M_{i}\) 上的自同态，以及 M 的一个把 \(M_{i}\) 映到 \(M_{i+1}\) 的一个子模上的自同态）。
\item
  反之，设 N 是 M 的一个子模，它是子模 \(N_{i} \subset M_{i}\) 的直和；假设诸零化理想 \(A\beta_{i}\) （对应诸 \(N_{i}\) ， \(1 \leqslant i \leqslant n\) ）满足上述条件。证明 N 是 M 的一个特征子模。由此推出，若 M 的两个特征子模具有相同的不变因子，则它们相等。
\end{enumerate}

\begin{enumerate}
\def\labelenumi{\arabic{enumi})}
\setcounter{enumi}{16}
\tightlist
\item
  设 A 为主理想整环，\(\alpha\) 是 A 的一个非零元素，并设 E 为模 A/A \(\alpha\) 。则乘积 \(E^{n}\) 同构于 \(A^{n}/(\alpha A^{n})\) 。
\end{enumerate}

\begin{enumerate}
\def\labelenumi{\alph{enumi})}
\item
  设 \(u\) 是从 \(\mathbf{E}^n\) 到 \(\mathbf{E}^m\) 的一个线性映射；证明 \(u\) 可由一个线性映射 \(v\) （从 \(\mathbf{A}^n\) 到 \(\mathbf{A}^m\) ）通过过渡到商而得到。
\item
  设 \(A\beta_{i}\) ( \(1 \leqslant i \leqslant p \leqslant \inf(m, n)\) ) 是 \(v(A^{n})\) 关于 \(A^{m}\) 的不变因子，按递减顺序排列；设 \(q \leqslant p\) 为使得 \(\alpha\) 不整除 \(\beta_{i}\) 的指标 i 中最大者，并设 \(\delta_{i}\) 为 \(\alpha\) 与 \(\beta_{i}\) 的一个最大公因子，对 \(1 \leqslant i \leqslant q\) 。证明 u 的核 \(u^{-1}(0)\) 是 n - q 个同构于 E 的模与 q 个分别同构于 \(A/A\delta_{i}\) （ \(1 \leqslant i \leqslant q\) ）的循环模的直和。若 \(\alpha = \gamma_{i}\delta_{i}\) ( \(1 \leqslant i \leqslant q\) )，证明模 \(u(E^{n})\) 具有不变因子 \(A\gamma_{i}\) 。
\end{enumerate}

\begin{enumerate}
\def\labelenumi{\arabic{enumi})}
\setcounter{enumi}{17}
\tightlist
\item
  设 \(\mathbf{C}\) 为一个主理想整环。考虑线性方程组
\end{enumerate}

\[
\sum_ {j = 1} ^ {n} \alpha_ {i j} \xi_ {j} = \beta_ {i} (1 \leqslant i \leqslant m)
\]

其中系数与右端项都属于 C，设 A 表示 m 行 n 列的矩阵 \((\alpha_{ij})\) ，而 B 表示通过添加第 \((n+1)\) 列 \((\beta_{i})\) 扩充 A 而得到的矩阵。证明该方程组在 \(C^{n}\) 中至少有一个解，当且仅当下列条件满足：1) A 与 B 具有相同的秩 p；2) A 的 p 阶子式的一个最大公因子等于 B 的 p 阶子式的一个最大公因子（利用 VII 第 20 页命题 4 以及习题 9，c））。

\begin{enumerate}
\def\labelenumi{\arabic{enumi})}
\setcounter{enumi}{18}
\tightlist
\item
  设 \(\mathbf{C}\) 是一个主理想整环。考虑由 \(m\) 个«线性同余式»组成的方程组 \(\sum_{j=1}^{n} \alpha_{ij} \xi_j \equiv \beta_i (\bmod \alpha)\) ，其中系数、右端项以及元素 \(\alpha\) 都属于 C，且 \(\alpha\) 非零；设 \(A\) 表示矩阵 \((\alpha_{ij})\) ，设 \(B\) 表示通过添加列 \((\beta_i)\) 扩充 A 而得到的矩阵，并设 \(\delta_k\) （相应地 \(\delta_k'\) ) 为 \(k\) 阶子式（\(A\) 的，相应地 \(B\) 的）的一个最大公因子。则该方程组在 \(\mathbf{C}^n\) 中至少有一个解，当且仅当：\(a)\) 若 \(m \leqslant n\) ，\(\alpha^m\) , \(\alpha^{m-1}\delta_1\) , \ldots, \(\alpha\delta_{m-1}\) , \(\delta_m\) 的一个最大公因子等于 \(\alpha^m\) , \(\alpha^{m-1}\delta_1'\) , \ldots, \(\alpha\delta_{m-1}'\) , \(\delta_m'\) 的一个最大公因子；
\end{enumerate}

\begin{enumerate}
\def\labelenumi{\alph{enumi})}
\setcounter{enumi}{1}
\tightlist
\item
  若 \(m > n\) ，\(\alpha^{n+1}\) , \(\alpha^n \delta_1, \ldots, \alpha \delta_n\) 的一个最大公因子等于 \(\alpha^{n+1}\) , \(\alpha^n \delta_1', \ldots, \delta_{n+1}'\) 的一个最大公因子。
\end{enumerate}

将线性同余式方程组化为线性方程组，并利用习题 18。

\begin{enumerate}
\def\labelenumi{\arabic{enumi})}
\setcounter{enumi}{19}
\tightlist
\item
  \begin{enumerate}
  \def\labelenumii{\alph{enumii})}
  \tightlist
  \item
    设 \(\mathbf{C}\) 是一个主理想整环，并且设
  \end{enumerate}
\end{enumerate}

\[
f (x _ {1},.., x _ {p}) = \sum_ {(i _ {k})} \alpha_ {i _ {1} \dots i _ {p}} \xi_ {1, i _ {1}} \dots \xi_ {p, i _ {p}}
\]

是一个 \(p\) -线性形式，定义在 \(\mathbf{E}^p\) 上，其中 \(\mathbf{E}\) 是 C-模 \(\mathbf{C}^n\) （而 \(\xi_{ij}\) 表示第 \(j\) 个坐标，对应于 \(x_i\) ）。证明若 \(\delta\) 是诸系数 \(\alpha_{i_1\dots i_p}\) 的一个最大公因子，则存在一个元素 \((a_k)\in \mathbf{E}^p\) ，使得 \(f(a_1,\dots,a_p) = \delta\) （化归到 \(\delta = 1\) 的情形，并对 \(p\) 作归纳，利用 VII，第 18 页，定理 1）。

\begin{enumerate}
\def\labelenumi{\alph{enumi})}
\setcounter{enumi}{1}
\tightlist
\item
  设 n \textgreater{} 1 为一个整数，且 \(\alpha_{i}\) ( \(1 \leqslant i \leqslant n\) ) 为 C 中以 \(\delta\) 为一个最大公因子的 n 个元素。证明存在 C 上一个 n 阶方阵 A，其第一列为 \((\alpha_{i})\) 且其行列式为 \(\delta\) （利用 a））。
\end{enumerate}

\begin{enumerate}
\def\labelenumi{\arabic{enumi})}
\setcounter{enumi}{20}
\tightlist
\item
  设 \(C\) 是一个主理想整环，并设 \(A = (\alpha_{ij})\) 为一个 \(n\) 阶方阵，定义在 \(C\) 上。
\end{enumerate}

\begin{enumerate}
\def\labelenumi{\alph{enumi})}
\item
  证明存在 C 上一个 n 阶方阵 U，其行列式为 1，使得矩阵 UA 的最后一列的前 \((n-1)\) 个元素为零，且最后一个元素等于 A 的第 n 列诸元素的一个最大公因子（利用习题 20，b））。
\item
  由 a) 推出，存在 C 上一个 n 阶方阵 V，其行列式为 1，使得矩阵 \(VA = (\beta_{ij})\) 是下三角的（«A 的埃尔米特简化型»）。
\item
  若 W 是 C 上一个 n 阶方阵，其行列式为 1，且 \(WA = (\gamma_{ij})\) 是下三角的，证明 \(\beta_{ii}\) 与 \(\gamma_{ii}\) 相伴，对 \(1 \leqslant i \leqslant n\) 。
\end{enumerate}

\begin{enumerate}
\def\labelenumi{\arabic{enumi})}
\setcounter{enumi}{21}
\tightlist
\item
  一个主理想整环上的模M被称为具有有限秩，如果存在一个整数 \(m \geqslant 0\) 具有如下性质：M 的每个有限生成子模都包含在一个由至多 m 个元素生成的子模中。具有该性质的最小整数 m 称为 M 的秩。
\end{enumerate}

\begin{enumerate}
\def\labelenumi{\alph{enumi})}
\item
  证明若 \(\mathbf{M}\) 具有有限秩，则每个商模都具有有限秩，且至多等于 \(\mathbf{M}\) 的秩。
\item
  对于无挠模 M，上述定义的秩概念与 II，第 315 页中定义的概念一致。
\item
  证明若 \(\mathbf{M}\) 具有有限秩，则 \(\mathbf{M}\) 的每个子模都具有有限秩，且至多等于 \(\mathbf{M}\) 的秩。
\item
  设 \(\mathbf{M}\) 是一个 \(\pi\) -准素模，满足 \(\mathbf{M}(\pi) = \mathbf{M}\) （采用 VII，第 7 页的记号）。证明若 \(\mathbf{M}\) 具有有限秩，则它是有限生成的。
\item
  设 M 为一个没有无限高元素的 \(\pi\) -准素模。证明若 M 具有有限秩，则它是有限生成的，并且秩的概念与
\end{enumerate}

习题 8 中所定义的一致。（注意若 \(\mathbf{M}(\pi)\) 是有限生成的，则 \(\mathbf{M}\) 也是有限生成的，这可通过证明 \(\mathbf{M}(\pi)\) 的元素在 \(\mathbf{M}\) 中的高度是有界的来得到；利用 VII 第 55 页习题 5 证明这一点。）\(f)\) 设 \(\mathbf{M}\) 是一个 \(\pi\) -准素模，具有有限秩 \(r\) ，并设 \(\mathbf{M}_0\) 表示 \(\mathbf{M}\) 中由所有无限高元素组成的子模。证明 \(\mathbf{M}_0\) 是一个可除模，并且是有限个（\(p \leqslant r\) 个）不可分解子模 \(\mathrm{U}_{\pi}\) （VII，第 54 页，习题 3）的直和，且 \(\mathbf{M}\) 是 \(\mathbf{M}_0\) 与一个有限生成子模 \(\mathbf{N}\) （秩为 \(r - p\) ）的直和。（注意 \(\mathbf{M} / \mathbf{M}_0\) 是有限生成的，利用 \(e\) )；然后应用 VII 第 54 页习题 3。）\(g)\) 由上述内容推出，一个 A-模 \(\mathbf{M}\) 具有有限秩 \(r\) ，当且仅当商模 \(\mathbf{M} / \mathbf{T}\) （\(\mathbf{M}\) 对其挠子模 \(\mathbf{T}\) 的商）具有有限秩 \(p \leqslant r\) ，且每个 \(\pi\) -准素分量（属于 \(\mathbf{T}\) ）的秩都 \(\leqslant r - p\) ，并且其中至少有一个的秩为 \(r - p\) 。

\begin{enumerate}
\def\labelenumi{\alph{enumi})}
\setcounter{enumi}{7}
\item
  设 \((p_{n})\) 是互异素数的序列；设 E 为循环 Z-模 \(Z/p_{n}^{2}Z\) 的直和，并设 \(e_{n}\) 表示这样的元素：其坐标除第 n 个外均为零，而第 n 个坐标等于 \(1 \mod p_{n}^{2}\) 的剩余类。设 M 为 Z-模 \(E \times Q\) 的由元素 \((e_{n}, 1/p_{n})\) 生成的子模。证明 M 是秩为 2 的 Z-模，M 的挠子模 T 由元素 \((p_{n}e_{n}, 0)\) 生成，且 T 不是 M 的直因子。（若 \(\bar{x}_{n}\) 是 \((e_{n}, 1/p_{n}) \mod T\) 的类，证明对所有 \(m \neq n\) 存在类 \(\bar{y}_{nm} \mod T\) 使得 \(\bar{x}_{n} = p_{m}\bar{y}_{nm}\) ，但不存在元素 \(x_{n}\) 属于类 \(\bar{x}_{n}\) ，使得对所有 \(m \neq n\) 存在 \(y_{nm} \in M\) 满足 \(x_{n} = p_{m}y_{nm}\) 。）
\item
  若 \(\mathbf{M}\) 是有限生成的 A-模，证明 \(\mathbf{M}\) 的秩等于 \(\gamma (\mathbf{M})\) ，即 \(\mathbf{M}\) 的生成集的基数中最小者。
\item
  设 M 是一个有限秩的挠 A-模。证明：若 N 是 M 的真子模，则 N 不能与 M 同构（利用 f）。
\end{enumerate}

\begin{enumerate}
\def\labelenumi{\arabic{enumi})}
\setcounter{enumi}{22}
\item
  设 A 为主理想整环，且 M 为有限生成的 A-模。证明：若 N 和 P 是两个 A-模，使得 \(M \oplus N\) 与 \(M \oplus P\) 同构，则 N 与 P 同构。（化归到 M 为循环模的情形；把 \(M \oplus N\) 与 \(M \oplus P\) 等同起来，我们可以考虑一个以两种方式分解为子模直和 \(F \oplus N\) 与 \(G \oplus P\) 的 A-模 E，其中 F 和 G 是循环的且同构的。在 F 和 G 同构于 A 的情形，证明若 \(N \neq P\) ，则 N 是 \(N \cap P\) 与一个同构于 A 的子模的直和。若 F 和 G 同构于 \(A/A\pi^{n}\) ，其中 \(\pi\) 是 A 的一个不可约元，考虑 F、G 的生成元 a、b，并考察 \(\pi^{n-1}a \notin G\) 、\(\pi^{n-1}b \notin F\) 的情形，最后是 \(\pi^{n-1}a \in G\) 与 \(\pi^{n-1}b \in F\) 同时成立的情形；在最后这种情形，考虑 E 的元素 \(a + b\) 及其生成的子模。）
\item
  阿贝尔群 G 同构于某个交换域 K 的非零元素构成的乘法群 \(K^{*}\) 的一个子群，当且仅当 G 的挠子群 T 的秩 \(\leqslant 1\) （习题 22）。（关于此条件的必要性，参见 V，第 79 页，命题 2；利用同一参考文献证明该条件对 T 同构于某个域 E 的 \(E^{*}\) 的一个子群是充分的。对所有 \(\alpha \in G/T\) ，令 \(g_{\alpha}\) 为 \(\alpha\) 在 G 中的一个陪集代表元，那么，把 T 与 \(E^{*}\) 的一个子群等同起来，我们可写成 \(g_{\alpha}g_{\beta} = e_{\alpha\beta}g_{(\alpha\beta)}\) ，对所有 \(\alpha\) , \(\beta \in G/T\) ，其中 \(e_{\alpha\beta} \in E\) 。证明存在一个以 \((b_{\alpha})_{\alpha \in G/T}\) 为基的 E-代数 B，使得 \(b_{\alpha}b_{\beta} = e_{\alpha\beta}b_{(\alpha\beta)}\) 对所有 \(\alpha\) , \(\beta\) 成立，并利用 VI，第 33 页，习题 20，c) 证明 B 是一个整环。则 B 的分式域 K 满足所需的条件。）
\item
  将 VII，第 18 页的定理 1 推广到环 A 上的左模，其中 A 不一定是交换的，但 A 的每个左理想和每个右理想都是主理想（这样的环的一个例子是张量积 \(\mathbf{K}[\mathbf{X}] \otimes_{\mathbf{K}} \mathbf{D}\) ，其中 \(\mathbf{K}\) 是一个交换域，且 \(\mathbf{D}\) 是一个 K-代数，它是一个域但不一定是交换的）。
\end{enumerate}

\BourbakiExerciseGroup

本节中的所有域，除非另有明确说明，均为交换域。

\begin{enumerate}
\def\labelenumi{\arabic{enumi})}
\item
  设 u 是有限维向量空间 E 的一个自同态，并设 V 是 E 中在 u 下不变的子空间。证明 u 在 V 上的限制的特征多项式整除 u 的特征多项式。由此得出 Cayley-Hamilton 定理的一个初等证明（化归到 \(E_{u}\) 是循环的情形，并利用 VII 第 30 页的公式 (2)）。
\item
  \begin{enumerate}
  \def\labelenumii{\alph{enumii})}
  \tightlist
  \item
    证明域 K 上的每个方阵都相似于其转置（化归到 Jordan 矩阵的情形）。
  \end{enumerate}
\end{enumerate}

\begin{enumerate}
\def\labelenumi{\alph{enumi})}
\setcounter{enumi}{1}
\item
  证明在有理整数环 \(\mathbf{Z}\) 上，矩阵 \(\begin{pmatrix} 8 & 2 \\ 0 & 1 \end{pmatrix}\) 不相似于其转置。
\item
  设 \(u\) 是域 \(\mathbf{K}\) 上有限维向量空间的一个自同态。则每个特征值 \(\lambda\) （属于 \(u\) ）也是转置自同态 \(^t u\) （属于 \(E^*\) ）的特征值。此外，若 \(\lambda\) 与 \(\mu\) 是 \(u\) 的两个不同的特征值，\(x\) 是 \(u\) 的对应于 \(\lambda\) 的特征向量，且 \(y^*\) 是 \(^t u\) 的对应于 \(\mu\) 的特征向量，则 \(\langle x, y^* \rangle = 0\) 。
\end{enumerate}

\begin{enumerate}
\def\labelenumi{\arabic{enumi})}
\setcounter{enumi}{2}
\item
  设 u 是有限维向量空间 E 的一个自同态，令 \(\lambda\) 为 u 的一个特征值，并令 v 为自同态 \(u - \lambda\) 。证明 E 是 \(v^{-1}(0)\) 与 \(v(E)\) 的直和，当且仅当 \(\lambda\) 是 u 的极小多项式的单根。
\item
  通过向基域添加一个超越元，证明 VII 第 36 页命题 10 中的公式 (6) 可以从公式 (8) 推出。通过将多项式 q 分解为线性因式来直接证明后者。
\item
  \begin{enumerate}
  \def\labelenumii{\alph{enumii})}
  \tightlist
  \item
    设 K 是一个代数闭域，A 是 K 上一个 n 阶方阵，呈对角分块 \(\text{diag}(A_i)\) （由诸下三角矩阵 \(A_i\) 组成）的形式，使得 \(A_i\) 的对角元都等于 K 的同一个元素 \(\alpha_i\) ( \(1 \leqslant i \leqslant r\) )，且 \(\alpha_i \neq \alpha_j\) 对 \(i \neq j\) 成立。证明每个与 A 可交换的方阵 B 都是对角分块 \(\text{diag}(B_i)\) ，其中 \(B_i\) 与 \(A_i\) 同阶，对 \(1 \leqslant i \leqslant r\) 。
  \end{enumerate}
\end{enumerate}

\begin{enumerate}
\def\labelenumi{\alph{enumi})}
\setcounter{enumi}{1}
\tightlist
\item
  设 M 是 K 上一些 n 阶可交换方阵组成的一个集合。证明通过对 M 中所有矩阵施加相同的相似变换，可使它们都成为下三角矩阵的分块对角矩阵 \(\text{diag}(X_i)\) ，其中 \(X_i\) 的个数 r 以及 \(m_i\) （即 \(X_i\) 对每个 i 的阶）对 M 中每个矩阵都相同，且 \(X_i\) 的所有特征值都相等。（利用 a) 化归到 r = 1 的情形，然后利用 VII 第 41 页引理 3 证明在此情形 M 中所有矩阵都有非零的公共特征向量，再对 n 作归纳完成证明。）
\end{enumerate}

\begin{enumerate}
\def\labelenumi{\arabic{enumi})}
\setcounter{enumi}{5}
\item
  交换域 K 上两个由 n 阶方阵组成的对 \((A_{1}, A_{2})\) 与 \((B_{1}, B_{2})\) 称为等价的，如果存在 K 上两个可逆方阵 P 和 Q，使得 \(A_{1} = PB_{1}Q\) 且 \(A_{2} = PB_{2}Q\) 。若 \(A_{1}\) 与 \(B_{1}\) 可逆，证明这两个对等价当且仅当矩阵 \(XA_{1} + A_{2}\) 与 \(XB_{1} + B_{2}\) 在环 K{[}X{]} 上等价。
\item
  设 \(A\) 是一个 \(m\) 行 \(n\) 列、秩为 \(r\) 的矩阵，定义在域 \(K\) 上。若 \(P\) 是一个 \(m\) 阶方阵而 \(Q\) 是一个 \(n\) 阶方阵（均在 \(K\) 上），使得 \(PAQ = A\) ，证明存在一个矩阵 \(P'\) 与 \(P\) 相似，以及一个矩阵 \(Q'\) 与 \(Q\) 相似，使得
\end{enumerate}

\[
P ^ {\prime} = \left( \begin{array}{c c} R & S \\ 0 & U \end{array} \right), \quad Q ^ {\prime} = \left( \begin{array}{c c} R ^ {- 1} & 0 \\ S ^ {\prime} & U ^ {\prime} \end{array} \right)
\]

其中R是一个r阶可逆方阵，U（相应地 \(U'\) ) 是 m - r 阶（相应地，n - r 阶）方阵，而 S（相应地， \(S'\) ）是一个具有 r 行和 m - r 列（相应地，n - r 行和 m 列）的矩阵。

\begin{enumerate}
\def\labelenumi{\arabic{enumi})}
\setcounter{enumi}{7}
\item
  设 E 是域 K 上有限维为 n 的向量空间，u 是 E 的一个自同态，V 是 E 中在 u 下不变的子空间，\(u|V\) 是 u 在 V 上的限制。若 \(f_{i}\) ( \(1 \leqslant i \leqslant n\) )（相应地 \(g_{j}\) ( \(1 \leqslant j \leqslant m\) )）是 u（相应地 \(u|V\) ）的相似不变量，证明 \(g_{j}\) 整除 \(f_{j+n-m}\) ，对 \(1 \leqslant j \leqslant m\) （参见 VII，第 64 页，习题 8，d)）。陈述并证明 u 诱导的 E/V 的自同态的类似性质。
\item
  设 \(A\) 是一个 \(n\) 阶方阵，\(B\) 是一个 \(m\) 阶方阵，\(C\) 是一个 \(m\) 行 \(n\) 列的矩阵，定义在域 \(K\) 上。把方程 \(X A = B X + C\) （其中 \(X\) 是一个 \(m\) 行 \(n\) 列的未知矩阵，定义在 \(K\) 上）化为多项式环 \(K[Y]\) 中的线性同余式方程组（对所有 \(f \in K[Y]\) ，注意 \(X f(A) = f(B)X + H(f)\) 对某个完全确定的矩阵 \(H(f)\) （ \(m\) 行 \(n\) 列）成立；令 \(\oplus_{i} E_i\) 为模 \(E\) （与 \(A\) 相伴，见第 1 节）分解为
\end{enumerate}

循环模的直和分解，这些循环模的零化理想即是 A 的异于 1 的相似不变量 \(f_{i}\) （命题 2），并令 \(a_{i}\) 为诸 \(E_{i}\) 的生成元；同样地定义子模 \(F_{j}\) （由 \(b_{j}\) 生成）以及对应于 B 的多项式 \(g_{j}\) ；写出 \(Xf_{i}(A)a_{i}=0\) ；然后把 \(Xa_{i}\) 与 \(H(f_{i})a_{i}\) 按诸 \(F_{j}\) 分解，并证明这样从方程 \(Xf_{i}(A)a_{i}=0\) 得到的条件对解的存在性是必要且充分的）。

\begin{enumerate}
\def\labelenumi{\arabic{enumi})}
\setcounter{enumi}{9}
\tightlist
\item
  \begin{enumerate}
  \def\labelenumii{\alph{enumii})}
  \tightlist
  \item
    证明：在任意域 K 上，矩阵
  \end{enumerate}
\end{enumerate}

\[
\left( \begin{array}{c c c c c c} \lambda & \alpha & \beta_ {13} & \beta_ {14} & ... & \beta_ {1n} \\ 0 & \lambda & \alpha & \beta_ {24} & ... & \beta_ {2n} \\ 0 & 0 & \lambda & \alpha & ... & \beta_ {3n} \\ \cdots & \cdots & \cdots & \cdots & \cdots & \cdots \\ 0 & 0 & 0 & 0 & ... & \alpha \\ 0 & 0 & 0 & 0 & ... & \lambda \end{array} \right)
\]

相似于若尔当矩阵 \(U_{n,\lambda}\) ，只要 \(\alpha \neq 0\) 。

\begin{enumerate}
\def\labelenumi{\alph{enumi})}
\setcounter{enumi}{1}
\item
  若 \(\lambda \neq 0\) 且 \(\mathbf{K}\) 的特征为 0，证明 \((U_{n,\lambda})^m\) 相似于若尔当矩阵 \(U_{n,\lambda^m}\) ，对每个整数 \(m\) （正的或负的）。
\item
  证明 \((U_{n,0})^m = 0\) 对 \(m \geqslant n\) 。对 \(0 < m < n\) ，设 \(n - 1 = km + q\) ，其中 \(0 \leqslant q < m\) ；证明 \((U_{n,0})^m\) 有 \(q + 1\) 个等于 \(\mathbf{X}^{k+1}\) 的相似不变量，以及 \(m - q - 1\) 个等于 \(\mathbf{X}^k\) 的相似不变量（把 \(\mathbf{K}^n\) 的典范基按适当的顺序排列）。
\item
  类似地，确定 \(f(U_{n,\lambda})\) （ \(f \in \mathbf{K}[\mathbf{X}]\) ）的相似不变量。
\item
  对特征为 \(p \neq 0\) 的域 K 讨论同样的问题。
\end{enumerate}

\begin{enumerate}
\def\labelenumi{\arabic{enumi})}
\setcounter{enumi}{10}
\item
  设 K 为一个代数闭域，设 A 为 K 上的一个可逆矩阵，且 m \textgreater{} 0 是一个整数。证明：如果 m 不是 K 的特征的倍数，则存在一个多项式 \(g(\mathbf{X}) \in \mathbf{K}[\mathbf{X}]\) 使得 \((g(A))^{m} = A\) （利用凯莱-哈密顿定理及 VII，第 54 页，习题 28）。把方程推广到 \(f(U) = A\) ( \(f \in K[X]\) ，U 为 K 上未知的 \(n \times n\) 矩阵）。
\item
  \begin{enumerate}
  \def\labelenumii{\alph{enumii})}
  \tightlist
  \item
    设 \(A\) 与 \(B\) 是两个 \(n\) 阶方阵，定义在交换域 \(\mathbf{K}\) 上，并设 \(f_i\) 与 \(g_i\) ( \(1 \leqslant i \leqslant n\) ) 分别是它们的相似不变量。证明在 \(\mathbf{K}\) 上由方阵 \(U\) （ \(n\) 阶且满足 \(UA = BU\) ）组成的向量空间同构于如下定义的空间 \(\mathbf{M}_0 = \mathbf{M} / \mathbf{N}\) ：\(\mathbf{M}\) 是方阵 \((u_{ij}(\mathbf{X}))\) 组成的空间，这些方阵是 \(n\) 阶的、定义在 \(\mathbf{K}[\mathbf{X}]\) 上，且满足 \(u_{ij}(\mathbf{X}) f_j(\mathbf{X})\) 是 \(g_i(\mathbf{X})\) 的倍式，对每对指标 \(i, j\) ；而 \(\mathbf{N}\) 是 \(\mathbf{M}\) 的子空间，由满足 \(u_{ij}(\mathbf{X})\) 是 \(g_i(\mathbf{X})\) 的倍式（对每对指标 \(i, j\) ）的那些矩阵组成。（把 \(U, A\) 与 \(B\) 分别看作自同态 \(u, v\) 与 \(w\) （属于 \(E = K^n\) ）的矩阵，并注意 \(u\) 是一个 \(K[X]\) -线性映射，从 \(E_v\) 到 \(E_w\) ；然后把 \(E_v\) 与 \(E_w\) 表示为商模。）
  \end{enumerate}
\end{enumerate}

\begin{enumerate}
\def\labelenumi{\alph{enumi})}
\setcounter{enumi}{1}
\item
  当 B = A 时，满足 UA = AU 的矩阵 U 构成一个环 P；此时 M 是一个方阵环，N 是 M 的一个双边理想，而 P 同构于商环 M/N。证明 P 与环 K{[}A{]} （由 \(I_{n}\) 和 A 生成）相同，当且仅当 A 的极小多项式等于它的特征多项式。
\item
  对任意的 \(A\) 与 \(B\) ，证明 \(\mathbf{M}_0\) 的维数等于 \(\sum_{i,j} m_{ij}\) ，其中
\end{enumerate}

\(m_{ij}\) 是 \(f_i\) 与 \(g_j\) 的最大公因子的次数（与 \(a\) 的方法相同）。

\begin{enumerate}
\def\labelenumi{\alph{enumi})}
\setcounter{enumi}{3}
\tightlist
\item
  证明矩阵 \(U \in \mathbf{M}_0\) 的秩的最大值等于 \(\sum_{k} d_k\) ，其中
\end{enumerate}

\(d_{k}\) 是 \(f_{k}\) 与 \(\pmb{g}_{k}\) 的最大公因子的次数（利用习题 8）。

\begin{enumerate}
\def\labelenumi{\alph{enumi})}
\setcounter{enumi}{4}
\tightlist
\item
  把 a)、c) 和 d) 的结果推广到如下情形：A 是 n 阶方阵，B 是 m 阶方阵，而 U 是 m 行 n 列的矩阵。
\end{enumerate}

\begin{enumerate}
\def\labelenumi{\arabic{enumi})}
\setcounter{enumi}{12}
\tightlist
\item
  设 \((\mathbf{Z}_{ij})\) ( \(1 \leqslant i, j \leqslant n\) ) 是 \(n^2\) 个不定元，定义在域 \(\mathbf{K}_0\) 上。设 \(\mathbf{A}\) 为多项式环 \(\mathbf{K}_0[\mathbf{Z}_{11}, \ldots, \mathbf{Z}_{nn}]\) ，并设 \(\mathbf{K} = \mathbf{K}_0(\mathbf{Z}_{ij})\) 为其分式域；令 \(\chi_U\) 为矩阵 \(U = (\mathbf{Z}_{ij})\) 在 \(\mathbf{K}\) 上的特征多项式。
\end{enumerate}

\begin{enumerate}
\def\labelenumi{\alph{enumi})}
\item
  证明 \(\chi_{U}\) 是一个可分的不可约多项式（证明若 \(\chi_{U}\) 在 K 上可约，则它将是 \(K_{0}[X, Z_{ij}]\) 中两个关于 X 的次数 \textgreater0 的多项式之积；由此推出 K 上每个 n 阶方阵的特征多项式都会是可约的；然后从系数在 \(K_{0}\) 上代数无关的多项式的例子导出矛盾；用类似的方法证明可分性）。
\item
  证明 U 相似于一个每个元素都等于 0、1 或多项式 \(\chi_{U}\) 的某个系数的矩阵（参见 VII，第 30 页，公式 (2)）。由此推出：若 \(K_{0}\) 是一个无限域，p 是 A 中一个多项式，使得 \(p(u_{ij}^{\prime}) = p(u_{ij}^{\prime\prime})\) 对每对相似的矩阵 \((u_{ij}^{\prime})\) 与 \((u_{ij}^{\prime\prime})\) （在 \(K_{0}\) 上）成立，则 p 属于由 \(\chi_{U}\) 的系数生成的 A 的子环。对有理函数 \(r \in K = K_{0}(Z_{ij})\) 陈述并证明类似的结果。
\end{enumerate}

\begin{enumerate}
\def\labelenumi{\arabic{enumi})}
\setcounter{enumi}{13}
\tightlist
\item
  设 K 是特征为 p 的非完全域，设 a 是 \(K - K^{p}\) 的一个元素，并设 L 是域 \(\mathrm{K}(a^{1/p})\) 。考虑 K-向量空间 \(E = L \otimes_{K} L\) ；设 u（相应地 \(u'\) ）是 E 的 K-自同态，由在代数 E 中乘以 \(a^{1/p} \otimes 1\) （相应地 \(1 \otimes a^{1/p}\) ）给出。证明 u 与 \(u'\) 都是半单的，u 与 \(u'\) 可交换，且 \(u - u'\) 是非零幂零的。
\end{enumerate}

\BourbakiUnnumberedChapter{历史注记}{历史注记（第六、七章）}

（第六章和第七章）

（注：括号中的罗马数字指本文末尾的参考文献。）

初等算术运算，特别是分数的运算，不可避免地会引出许多关于整数之间整除性的经验观察。但无论是巴比伦人（尽管他们是代数学的专家），还是埃及人（尽管他们处理分数时技巧娴熟），似乎都未曾掌握支配这些性质的一般规则，而在这方面开创先河的是希腊人。他们的算术著作，其精妙阐述可见于欧几里得（I）的第七卷和第九卷，丝毫不逊色于他们在数学其他分支中最杰出的发现。两个整数的最大公约数的存在性，在第七卷的开头便通过被称为“欧几里得算法”* 的程序得以证明；它构成了后续所有发展的基础（素数的性质、最小公倍数的存在性与计算等），其论证方式与上文第六章第一节中的并无实质差异；而最辉煌的成就则是由两个卓越的定理构成，它们证明了素数的无穷性（第九卷命题20），并给出了构造偶完全数的方法（参见第七章第53页习题24；事实上，这一方法给出了所有偶完全数，正如欧拉后来所证明的那样）。唯有素数分解的存在性和唯一性未被以一般方式证明；然而，欧几里得确实明确证明了每个整数都能被某个素数整除（第七卷命题31），以及以下两个命题（第九卷命题13和14）：

“如果有任意多个数成连比例，首项为 1，公比恒定 {[}即等比数列{]}，且 1 后面的那个数是素数，那么最大的数将不能被任何数整除，除非这些数出现在该数列中”（换句话说，一个幂 \(p^n\) ——素数 \(p\) 的幂——只能被 \(p\) 的幂整除，而这些幂的指数 \(\leqslant n\)）。

“如果一个数是能被 {[}给定的{]} 素数整除的最小数，则它不能被任何其他素数整除，最初 {[}给定的{]} 整除它的素数除外”（换句话说，不同素数 \(p_{1}, \ldots, p_{k}\) 的乘积没有 \(p_{1}, \ldots, p_{k}\) 之外的素因子）。

因此看来，如果欧几里得没有陈述一般定理，那仅仅是因为缺乏适当的术语和记号来表示整数的任意次幂。*

尽管仔细研究使人觉得欧几里得的文本很可能由几个连续的层次组成，每一层对应于算术发展的一个阶段**，但似乎这一演变完全发生在公元前5世纪初至公元前4世纪中叶之间，人们不得不钦佩它所体现的技巧和逻辑自信：算术领域下一次可与之相比的进步要等到两千年之后才出现。

正是所谓的「不定」或「丢番图」问题构成了数论后续发展的根源。今天使用的「丢番图方程」一词在历史上并不完全恰当；它通常被理解为具有整数系数的多项式方程（或方程组），只寻求整数解：当方程是「确定的」，即只有有限多个（实数或复数）解时，这样的问题通常是不可能的；但另一方面，当未知数多于方程数时，它往往有解。现在，虽然丢番图确实似乎是第一个考虑「不定」问题的人，但他只在特殊情况下寻求整数解，通常满足于找到有理数中的一个解（II）。那是他通常能够通过代数计算解决的一类问题，其中

* 为支持这一论点，还可以指出，完全数定理的证明基本上只是唯一素因子分解定理的另一个特例。此外，所有证据表明，从这一时期起，将一个明确的数分解为素数已是众所周知并经常使用的；但在高斯于《算术研究》开头给出的证明之前，找不到分解定理的完整证明（（VIII）第I卷，第15页）。

** 参见 B. L. van der WAERDEN，《毕达哥拉斯学派的算术》，Math. Ann.，第CXX卷（1947-49），第127页。前一版本遗留段落的一个例子是第九卷第21至34命题，它们讨论被2整除的最基本性质，无疑可追溯到素数一般理论尚未发展的时期。此外，众所周知，偶数和奇数范畴在早期毕达哥拉斯学派形而上学的思辨中起了重要作用，因此自然倾向于将这一部分归于他们。

未知数的算术性质并不相关*；此外，整除性理论只起次要作用（“素数”一词仅出现一次（（II），第五卷，问题9，第1卷，第334-335页），而互素数的概念仅在涉及如下定理时才被引用：该定理断言，两个互素数的商能成为平方数，仅当它们各自都是平方数）**。

对不定方程整数解的研究，真正开始于中世纪早期的中国和印度数学家。前者似乎是由编制历法的实际问题引导到这类考虑的（其中确定各种天文现象循环的共同周期，恰恰构成一个线性的“丢番图”问题）；无论如何，他们（确定在公元4世纪至7世纪之间）提出了求解一次同余式组的法则（参见VI，第33页，习题25）。至于印度人，其数学在公元5世纪至13世纪间繁荣发展，他们不仅知道如何有条理地处理（通过应用欧几里得算法）任意多个未知数的线性丢番图方程组***，而且他们率先攻克并解决了二次问题，其中包括“费马方程”的某些特殊情形 \(Nx^{2} + 1 = y^{2}\) （（III），第II卷，第87-307页）。

此处不宜详述次数 \textgreater1 的丢番图方程理论的历史，该理论经由费马、欧拉、拉格朗日和高斯的工作，在 19 世纪导向了代数整数理论。正如我们已指出的（参见第II章和第III章的历史注记），线性方程组的研究在这一时期相当被忽视，因为它似乎不再呈现值得关注的问题：特别是，既没有寻求任意方程组的一般存在条件，也没有描述解集。尽管如此，埃尔米特在他于 19 世纪

* 虽然丢番图关于不定问题的工作总是归结为单一未知数的问题，即通过数值选取其他未知数，以使其最终方程有解，但使用这一方法的主要原因似乎是他的记号不允许他同时用多个未知数进行计算；无论如何，他在整个计算过程中都记录他所做的数值代入，并在需要时通过写出代入变量的相容性条件并先解这个辅助问题来修改它们。换言之，他处理这些代入的数值就像我们处理参数一样，以至于他实际所做的相当于求一个给定代数簇或其子簇的有理参数表示（参见（II bis））。

世纪的数论研究中，不得不使用关于线性丢番图方程的若干引理，特别是整数系数线性变换的一个“简化形式”（（XIII），第 164 页和第 265 页）；最后，在 1858 年赫格尔给出了秩等于方程个数的方程组的存在条件之后，H. J. 史密斯在 1861 年定义了整数矩阵的不变因子，并得到了如下一般定理：这样的矩阵可化为我们在 VII，第 22 页，推论 1 中给出的“标准形”（XVII）。

但与此同时，随着高斯引入这一概念（参见第一章、第二章和第三章的历史注释），以及它在数论后续发展中所起的重要作用，阿贝尔群的概念逐渐被精确化。在他于《算术研究》中对给定判别式的二次型类所构成的有限阿贝尔群所作的特别深入的研究中，高斯很快意识到其中某些群并非循环群：“在这种情况下，”他说，“一个基” {[}也就是说，一个生成元{]} 不足以满足要求，有必要取两个或更多个，通过乘法和合成*，可以产生所有类»（（VIII），第 I 卷，第 374-375 页）。不能确定高斯是否打算用这些话来描述群作为循环群的直积的分解；然而，在《算术研究》的同一篇文章中，他证明了群中存在一个元素，其阶是所有元素阶的最小公倍数——换言之，他得到了群的最大不变因子的存在性（（VIII），第 I 卷，第 373 页）；另一方面，直积的概念他是知道的，因为在一份写于1801年但生前未发表的手稿中，他勾勒了把有限阿贝尔群分解为 \(p\) -群的直积的一般证明 ** （（VIII），第 II 卷，第 226 页）。无论如何，在1868年，高斯著作集的编辑谢林，受到这些结果（特别是他刚刚发现的这份手稿）的启发，证明了（仍然针对二次型类群）一般的分解定理 (XVIII)，其方法——正如两年后克罗内克 (XX) 以抽象术语所重复的那样——本质上就是我们上面使用过的方法 (VII，第18页，定理1)。至于无挠阿贝尔群，我们已经说过（参见第II章和第III章的历史注记），由高斯、阿贝尔和雅可比发展的椭圆函数和阿贝尔积分理论，是如何逐渐引导人们关注其结构的；第一个也是最著名的无限群分解为循环群直和的例子，是1846年狄利克雷在他关于代数数域单位的论文中给出的 (XI)。但直到1879年，有限生成阿贝尔群理论与史密斯定理之间的联系才被弗罗贝尼乌斯和施蒂克尔贝格 ((XXIII)，第10节) 认识到并明确使用。

大约在同一时期，矩阵（实或复元素）的相似理论也趋于完成。线性变换特征值的概念明确出现在常系数线性微分方程组理论中，拉格朗日 (VIa) 将其应用于小振动理论，拉格朗日 (VIb) 和拉普拉斯 (VIIa) 将其应用于行星的“长期”摄动。它隐含在许多其他问题中，这些问题也在18世纪中叶左右被研究，例如寻找圆锥曲线或二次曲面的轴的问题（首先由欧拉 (Va) 解决），或者（也由欧拉 (Vb) 发展的）刚体主惯性轴的研究（由德·塞格纳于1755年发现）；我们现在知道，它（以更隐蔽的形式）也涉及偏微分方程理论的初期，特别是振动弦方程。但是（撇开最后这种情况不谈），在柯西 (X) 之前，这些问题之间的联系几乎没有被认识到。此外，由于其中大多数问题涉及对称矩阵的使用，特征值最初之所以被研究，主要是因为这些矩阵；我们将在本专著专门讨论埃尔米特算子的各章之后的历史注记中更详细地回到这一点；这里只需指出，早在1826年，柯西就证明了此类矩阵的特征值在相似变换下的不变性，并证明了它们对于 \(3 \times 3\) 对称矩阵 (Xa) 是实数，三年后 (Xb) 他将此结果推广到任意实对称矩阵。* 由莫比乌斯于1827年引入的投影的一般概念，迅速导致了对此类变换进行分类的问题（首先在2维和3维），这无非就是相应的矩阵在相似意义下的分类问题；但很长一段时间里，这个问题仅用19世纪中叶流行的“综合”方法处理，其（无论如何相当缓慢的）进展似乎并未以任何方式影响特征值理论。另一个几何问题，即圆锥曲线或二次曲面“束”的分类，情况则不同，从现代观点来看，这相当于研究矩阵 \(U + \lambda V\) ，其中 \(U\) 与 \(V\) 是两个对称矩阵；西尔维斯特在 1851 年着手解决这个问题时，肯定就是本着这种精神，仔细考察（为所考虑的束找出“标准形”）矩阵 \(U + \lambda V\) 的子式当为 \(\lambda\) 代入一个使行列式为零的值 (XIV) 时会发生什么。特征值理论的纯代数方面同时也在发展；因此，大约在 1850 年，几位作者（包括西尔维斯特本人）证明了 \(U^n\) 的特征值是 \(n\) 次幂，即 \(U\) 的特征值的幂，而凯莱在 1858 年，在他引入矩阵算术的那篇论文（XVI）中，对任意方阵*宣布了“凯莱-哈密顿定理”，尽管他只对 \(2 \times 2\) 与 \(3 \times 3\) 矩阵给出了直接证明。最终，魏尔斯特拉斯在 1868 年，利用西尔维斯特的方法，得到了“束” \(U + \lambda V\) 的“标准形”，此时 \(U\) 与 \(V\) 是方阵，不必对称，只需满足条件 \(\det(U + \lambda V)\) 不恒为零；他由此推导出任意（复）方阵的初等因子的定义，并证明了它们刻画了矩阵在相似意义下的等价类（XIX）；顺便说一句，这些结果后来由若尔当在两年后部分地（且显然是独立地）重新得到**（XXI）。这里又是弗罗贝尼乌斯在1879年指出，魏尔斯特拉斯定理可以容易地从史密斯定理（推广到多项式的情形）推出（（XXII），§ 13）；他的方法正是我们上文给出的该定理证明的基础（VII，第35页）。

我们刚才提到了单变量多项式的整除性理论；多项式的除法问题自然必定在代数学的最早期就已出现，因为它是乘法的逆运算（后者甚至丢番图就已经知道，至少对于低次多项式是如此）；但可以想象，在变量的各次幂的连贯记法建立之前，几乎不可能以一般方式处理这个问题。事实上，在我们所知的16世纪中叶之前，几乎找不到多少“欧几里得”式多项式除法程序的例子***；而S.斯蒂文（基本上使用指数记法）似乎是第一个想到将“欧几里得算法”推广以求出两个多项式的最大公因式的人（（IV，第I卷，第54-56页））。除此之外，直到18世纪中叶，整除性理论一直局限于有理整数。正是欧拉在1770年开创了算术的新篇章，他多少有些轻率地将整除性概念推广到二次扩域中的整数：在寻求确定形如 \(x^{2} + cy^{2}\) ( \(x, y\) 与 \(c\) 为有理整数）的数的因子时，他令

\[
x + y \sqrt {- c} = (p + q \sqrt {- c}) (r + s \sqrt {- c}) (p, q, r, s \text { rational   integers })
\]

并且，对两边取范数，他毫不犹豫地断言， \(x^{2} + cy^{2}\) 的所有因子都可以这样得到，形如 \(p^{2} + cq^{2}\) （Vc）。换言之，

欧拉仿佛认为环 \(Z[\sqrt{-c}]\) 是主理想整环；稍后他又用类似的论证将“无穷递降”方法应用于方程 \(x^{3} + y^{3} = z^{3}\) （他将问题归结为求 \(p^{2} + 3q^{2}\) 的立方根，为此他令 \(p + q\sqrt{-3} = (r + s\sqrt{-3})^{3}\) ）。但拉格朗日在1773年就证明了（VIc），形如 \(x^{2} + cy^{2}\) 的数的因子并非都具有这种形式，这是后来在高斯及其后继者对分圆域中整除性的研究中将更加清楚地显现出来的根本困难的第一个例子*；一般来说，不可能将有理整数整除性的基本性质（如最大公因式的存在性和素因子分解的唯一性）直接推广到这些域。这里不适合描述库默尔如何针对分圆域（XII）**，以及随后戴德金和克罗内克如何针对任意代数数域，通过理想理论的发明成功地克服了这一巨大障碍——这是现代代数中最具决定性的进展之一。但戴德金始终对各种数学理论的基础充满好奇，他并不满足于这一成功；通过分析整除性关系的机制，他在一篇论文中奠定了格序群理论的基础（这篇论文当时不为同时代人所知，湮没无闻达30年之久），它无疑是公理化代数学最早的著作之一（XXIV）；除去记号上的差异，他的工作已非常接近我们在第六章§ 1中呈现的这一理论的现代形式。

从18世纪中叶起，当时的主要任务是寻找“代数基本定理”（参见第四、五章的历史注记）。这里我们无需提及达朗贝尔的尝试，他开创了使用无穷小演算的证明系列（参见《一般拓扑学》第八章的历史注记）。但欧拉在1749年以完全不同的方式处理了这个问题

* 高斯似乎曾一度希望 \(n\) 次单位根域中的整数环会是主理想整环；在他生前未发表的一份手稿中（（VIII），第II卷，第387-397页），他证明了三次单位根域中存在欧几里得除法过程，并给出五次单位根域中类似过程的一些提示；他利用这些结果，通过比欧拉的论证更正确的“无穷递降”方法证明了方程 \(x^3 + y^3 = z^3\) 在三次单位根域中无解，并指出可以将该方法推广到方程 \(x^5 + y^5 = z^5\) ，但在方程 \(x^7 + y^7 = z^7\) 处止步，说在这里无法先验地排除 \(x, y\) 与 \(z\) 不被 7 整除的情形。

(Vd)：对每个实系数多项式 f，他试图证明存在分解 \(f = f_{1}f_{2}\) ，分解为两个（非常数的）实系数多项式，这将使他通过对 f 的次数作归纳来证明“基本定理”。正如他所注意到的，甚至只需止步于第一个奇数次因子即可，因此所有困难都归结为考虑 f 的次数 n 为偶数的情况。欧拉随后把自己限制在所求因子次数各为 n/2 的情形，并指出通过适当的消元过程，可以把 \(f_{1}\) 与 \(f_{2}\) 的未知系数表示为一个实系数方程的根的有理函数，该方程的首末两项符号相反，因而至少有一个实根。但欧拉的证明只是一个草图，跳过了许多要点；直到 1772 年，拉格朗日才通过极其冗长而细致的分析成功解决了该证明提出的困难（VId），在此过程中他展示了在运用他新近创造的“伽罗瓦方法”方面的非凡技巧（参见第四章和第五章的历史注记）。

尽管如此，拉格朗日与欧拉及其同时代的所有人一样，毫不犹豫地在多项式的“根域”内进行形式推理（也就是说，用他的语言，考虑该多项式的“虚根”）；他那个时代的数学并未为这类论证提供任何合理性证明。高斯从一开始就坚决反对18世纪无限制的形式主义，在他的学位论文中强烈抨击了这种滥用（（VIII），第III卷，第3页）。但他不可能没有察觉到，这本质上是一个内在正确论证的表面上有缺陷的表述。几年后，我们还发现他（（VIII），第III卷，第33页；另见（VIIIBis））采用了欧拉论证的一个更简单的形式，该形式由丰瑟内于1759年提出（然而，丰瑟内未能利用它取得任何成果），以获得“基本定理”的一个新证明，在该证明中他小心地避免提及所有“虚”根：后者被巧妙地添加和特殊化不定元所取代。我们在正文中给出的本质上正是高斯的这个证明（VI，第26页，定理3），并利用代数扩张进行了简化。

因此，拓扑在“基本定理”中的作用被归结为这样一个定理：具有实系数的多项式在区间内不能在不取零值的情况下改变符号（博尔扎诺多项式定理）。该定理也是所有关于（具有实系数的）多项式实根分离判据的基础，这是19世纪最受青睐的代数课题之一*。在这项研究过程中，显而易见的是，起关键作用的是R的序结构，而非其拓扑结构*；例如，博尔扎诺多项式定理在全体实代数数域中仍然成立。这一思路最终导致了由E.阿廷和O.施赖埃尔（XXV）创建的序域的抽象理论；其最显著的结果之一无疑是发现了域上序关系的存在与该域的纯代数性质相关。这就是第六章第2节所呈现的理论。

\BourbakiUnnumberedChapter{参考文献}{参考文献}

\begin{enumerate}
\def\labelenumi{(\Roman{enumi})}
\tightlist
\item
  《欧几里得原本》，5卷，J. L. 海伯格编，莱比锡（托伊布纳），1883-88年。
\end{enumerate}

（I bis）T. L. 希思，《欧几里得几何原本》十三卷\ldots，3卷，剑桥，1908年。

\begin{enumerate}
\def\labelenumi{(\Roman{enumi})}
\setcounter{enumi}{1}
\tightlist
\item
  《丢番图全集》\ldots，2卷，P. 坦纳里编，莱比锡（托伊布纳），1893-95年。
\end{enumerate}

（II bis）T. L. 希思，《亚历山大的丢番图》，第2 \(^{nd}\) 版，剑桥，1910年。

\begin{enumerate}
\def\labelenumi{(\Roman{enumi})}
\setcounter{enumi}{2}
\item
  B. 达塔与A. N. 辛格，《印度数学史》，2卷，拉合尔（莫蒂拉尔·巴纳尔西·达斯出版社），1935-38年。
\item
  S. 斯蒂文，《数学著作》\ldots，A. 吉拉德编，莱顿（爱思唯尔出版社），1634 年，第 I 卷。
\item
  L. 欧拉：a)《无穷分析引论》（《全集》，(1) 系列，第IX卷，苏黎世-莱比锡-柏林（O. Füssli 与 B. G. Teubner 出版社），1945年，第384页）；b)《固体或刚体运动理论》（《全集》(2) 系列，第III卷，苏黎世-莱比锡-柏林（O. Füssli 与 B. G. Teubner 出版社），1948年，第200-201页）；c)《代数学完全入门》（《全集》(1) 系列，第I卷，莱比锡-柏林（Teubner 出版社），1911年，第422页）；d)《关于方程虚根的研究》（《全集》(1) 系列，第VI卷，莱比锡-柏林（Teubner 出版社），1921年，第78页）。
\item
  J.-L. 拉格朗日，《著作集》，巴黎（Gauthier-Villars 出版社），1867-1892年：a)《积分学中若干问题的解法》，第I卷，第520页；b)《关于交点运动长期方程的研究》，第VI卷，第655-666页；c)《算术研究》，第III卷，第695-795页；d)《关于方程虚根的形式》，第III卷，第479页；e)《不受任何加速力作用的任意物体旋转问题的新解法》，第III卷，第579-616页。
\item
  P. S. 拉普拉斯：a) 关于微分方程的特解及行星长期不等式的论文（《著作集》，第 VIII 卷，巴黎（戈蒂埃-维拉尔出版社），1891 年，第 325-366 页）；b) 关于行星和卫星长期不等式的论文（《著作集》，第 XI 卷，巴黎（戈蒂埃-维拉尔出版社），1895 年，第 49-92 页）。
\item
  C. F. 高斯，《著作集》，第 I 卷（哥廷根，1870 年），第 II 卷（同上，1876 年）及第 III 卷（同上，1876 年）。
\end{enumerate}

(VIII bis) 高斯关于整代数函数分解为一次或二次实因子的四个证明（奥斯特瓦尔德经典著作，第14号，莱比锡（托伊布纳），1904年）。

\begin{enumerate}
\def\labelenumi{(\Roman{enumi})}
\setcounter{enumi}{8}
\item
  N. H. 阿贝尔，《全集》，第 I 卷，西洛与李编，克里斯蒂安尼亚，1881 年。
\item
  A. L. 柯西：a)《无穷小计算在几何中的应用讲义》（《全集》第二辑，第五卷，巴黎（戈蒂埃-维拉尔出版社），1903年，第248页）；b)《关于用以确定行星长期不等式的方程》（《全集》第二辑，第九卷，巴黎（戈蒂埃-维拉尔出版社），1891年，第174页）。
\item
  P. G. 勒热纳-狄利克雷，《著作集》，第 I 卷，柏林（G. 赖默），1889 年，第 619-644 页。
\item
  E. 库默尔，《关于复数理论》，J. de Crelle，第 XLIII 卷（1847 年），第 319 页（《论文集》，第 I 卷，海德堡（Springer 出版社），1975 年，第 203 页）。
\item
  Ch. 埃尔米特，《著作集》，第 I 卷，巴黎（Gauthier-Villars 出版社），1905 年。
\item
  J. J. 西尔维斯特，《数学论文集》，第 I 卷，剑桥，1904 年：关于二阶直线与曲面接触的枚举，第 219 页（=《哲学杂志》，1851 年）。
\item
  W. R. 哈密顿，《四元数讲义》，都柏林，1853 年。
\item
  A. 凯莱，《数学论文集》，剑桥，1889-1898 年：关于矩阵理论的回忆录，第 II 卷，第 475-496 页（=《哲学汇刊》，1858 年）。
\item
  H. J. 史密斯，《数学论文集》，第 I 卷，牛津，1894 年；关于线性不定方程组与同余式组，第 367 页（《哲学汇刊》，1861 年）。
\item
  E. 谢林，《可复合算术形式的基本类》，哥廷根科学学会论文集，第 XIV 卷（1868-69 年），第 13 页。
\item
  K. 魏尔斯特拉斯，《数学著作》，第 II 卷，柏林（Mayer und Müller 出版社），1895 年：关于双线性与二次形式理论，第 19 页。
\item
  L. 克罗内克，《关于理想复数类数的若干性质的阐述》，柏林科学院月报（1870 年），第 881 页（=《著作集》，第 I 卷，莱比锡（Teubner 出版社），1895 年，第 273 页）。
\item
  C. 若尔当，《置换与代数方程论》。巴黎（Gauthier-Villars 出版社），1870 年，第 114-125 页。
\item
  G. 弗罗贝尼乌斯，《具有整数系数的线性形式理论》，《全集》，第 I 卷，海德堡（Springer 出版社），1968 年，第 482 页（= J. de Crelle，1879 年）。
\item
  G. 弗罗贝尼乌斯与 L. 施蒂克尔贝格，《关于可交换元素的群》，J. de Crelle，第 LXXXVI 卷（1879 年），第 217 页（= 弗罗贝尼乌斯，《全集》，第 I 卷，第 545 页）。
\item
  R. 戴德金，《数学全集》，第 II 卷，不伦瑞克（Vieweg 出版社），1932 年：关于数通过其最大公因子的分解，第 103 页。
\item
  \begin{enumerate}
  \def\labelenumii{\alph{enumii})}
  \tightlist
  \item
    E. 阿廷与 O. 施赖埃尔，《实域的代数构造》，汉堡大学数学讨论班论文集，第 V 卷（1927 年），第 83 页；b) E. 阿廷，《关于确定函数分解为平方和》（同上，第 100 页）；c) E. 阿廷与 O. 施赖埃尔，《实闭域的一种刻画》（同上，第 225 页）。
  \end{enumerate}
\end{enumerate}

\BourbakiUnnumberedChapter{记号索引}{记号索引}

\[
\mathbf {A} \left[ \left(\mathbf {X} _ {i}\right) _ {i \in I} \right], \mathbf {A} \left[ \mathbf {X} _ {1}, \dots , \mathbf {X} _ {n} \right], \mathbf {X} ^ {\nu}: \mathrm{IV}, \text { p. } 1.
\]

\[
\deg u: \mathrm{IV}, \text { p. } 3.
\]

\[
\mathbf {f} \circ \mathbf {g}: \mathrm{IV}, \text { p. } 4.
\]

\[
\mathrm{D} _ {i} \mathrm{P}, \mathrm{D} _ {\mathrm{X} _ {i}} \mathrm{P}, \frac {\partial \mathrm{P}}{\partial \mathrm{X} _ {i}}, \mathrm{P} _ {\mathrm{X} _ {i}} ^ {\prime}, \mathrm{D} ^ {\nu}, \mathrm{DP}, \frac {d \mathrm{P}}{d \mathrm{X}}, \mathrm{P} ^ {\prime}: \mathrm{IV}, \text { p. } 6.
\]

\[
\Delta^ {\nu}: \mathrm{IV}, \text { p. } 7.
\]

\[
\mathbf {K} ((\mathbf {X} _ {i}) _ {i \in I}), \deg r: \mathrm{IV}, \text { p. } 19 \text { and } 20.
\]

\[
\mathrm{D} _ {i} f, \mathrm{D} _ {\mathbf {X} _ {i}} f, \frac {\partial f}{\partial \mathbf {X} _ {i}}, f _ {\mathbf {X} _ {i}} ^ {\prime}, \mathrm{D} f, \frac {d f}{d \mathbf {X}}, f ^ {\prime}: \mathrm{IV}, \text {p.} 23.
\]

\[
\mathbf {A} \left[ \left[ \left(\mathbf {X} _ {i}\right) _ {i \in I} \right] \right], \mathbf {A} \left[ [ \mathrm{I} ] \right]: \mathrm{IV}, \text { p. } 24.
\]

\[
\omega : \mathrm{IV}, \text { p. } 25.
\]

\[
u (\mathbf {x}), u ((x _ {i}) _ {i \in I}), u (x _ {1}, \dots , x _ {n}): \mathrm{IV}, \text { p. } 29.
\]

\[
\Delta^ {\nu}, \mathrm{D} ^ {\nu}, \mathrm{D} _ {i}: \mathrm{IV}, \text { p. } 31 \text { and } 32.
\]

\[
\mathrm{D} _ {i} u, \mathrm{D} _ {\mathbf {X} _ {i}} u, \frac {\partial u}{\partial \mathbf {X} _ {i}}, u _ {\mathbf {X} _ {i}} ^ {\prime}, \mathrm{D} u, \frac {d u}{d X}: \mathrm{IV}, \mathrm{p.32}.
\]

\[
\mathrm{A} \{\mathrm{I} \}, \mathbf {f} \circ \mathbf {g}: \mathrm{IV}, \text { p. } 29 \text { and } 36.
\]

\[
\mathrm{T} _ {\mathbf {f}}: \mathrm{IV}, \text { p. } 36.
\]

\[
\exp \mathrm{X}, e ^ {\mathrm{X}}, e (\mathrm{X}), l (\mathrm{X}): \text { IV }, \text { p. } 39 \text { and } 40.
\]

\[
\exp f, \log g: \mathrm{IV}, \text { p. } 40.
\]

\[
\mathbf {M} ^ {\mathrm{H}}, \operatorname{Tr} _ {\mathrm{H} / \mathrm{G}}: \mathrm{IV}, \text { p. } 41 \text { and } 42.
\]

\[
\mathbf {T S} ^ {n} (\mathbf {M}), \mathbf {T S} (\mathbf {M}): \mathrm{IV}, \text { p. } 42.
\]

\[
\mathfrak {S} _ {p \mid q}, \mathfrak {S} _ {p, q}, \mathfrak {S} _ {p _ {1} \mid \dots \mid p _ {n}}: \mathrm{IV}, \text {   p.   } 43.
\]

\[
\gamma_ {k} (x), x \in \mathrm{M}: \mathrm{IV}, \text { p. } 45.
\]

\[
\varphi_ {\mathbf {M}}, \psi_ {\mathbf {M}}: \mathrm{IV}, \text { p. } 52.
\]

\[
\operatorname{Pol} _ {\mathbf {A}} ^ {q} (\mathbf {M}, \mathbf {N}), \operatorname{Pol} ^ {q} (\mathbf {M}, \mathbf {N}): \mathrm{IV}, \text { p. } 55.
\]

\[
\operatorname{Map} (\mathbf {M}, \mathbf {N}), \operatorname{Pol} _ {\mathbf {A}} (\mathbf {M}, \mathbf {N}), \operatorname{Pol} (\mathbf {M}, \mathbf {N}): \mathrm{IV}, \text { p. } 57.
\]

\[
s _ {k}, s _ {k, n}, \mathrm{A} [ \mathrm{X} _ {1}, \dots , \mathrm{X} _ {n} ] ^ {\text { sym }}: \mathrm{IV}, \text { p. } 61.
\]

\[
\mathrm{S} (\alpha), \mathrm{M} (\alpha): \mathrm{IV}, \text { p. } 65 \text { and } 66.
\]

\[
s _ {k}, \mathrm{A} [ [ \mathrm{X} ] ] ^ {\text { sym }}: \mathrm{IV}, \text { p. } 67 \text { and } 68.
\]

\[
\mathfrak {B} _ {k}: \mathrm{IV}, \text { p. } 70.
\]

\[
M (f, g, p, q), \operatorname{res} _ {p, q} (f, g), \operatorname{res} (f, g): \mathrm{IV}, \text {   p.   } 76.
\]

\[
\operatorname{dis} (f), f \text {   monic   polynomial }: \mathrm{IV}, \mathrm{p.~81}.
\]

\[
\operatorname{dis} _ {m} (f), f \text { polynomial   of   degree } \leqslant m: \mathrm{IV}, \text { p. } 83.
\]

\[
\Gamma (\mathrm{E}), \Gamma_ {p} (\mathrm{E}), \gamma_ {p}, \Gamma (h): \mathrm{IV}, \text { p. } 92, \text { ex. } 2.
\]

\[
\mathbf {Q}, \mathbf {F} _ {p}: \mathbf {V}, \text { p. } 1.
\]

\[
\mathbf {S} ^ {p ^ {f}}: \mathbf {V}, \mathrm{p.} 4.
\]

\(U_{m,\alpha}\) ：第七卷，第35页。

\[
\mathrm{K} [ \mathrm{S} ]: \mathrm{V}, \text { p. } 4.
\]

\[
\mathbf {S} ^ {p ^ {- f}}, \mathbf {A} ^ {p ^ {- \infty}}: \mathbf {V}, \text { p. } 5 \text { and } 6.
\]

\[
[ \mathbf {A}: \mathbf {K} ]: \mathbf {V}, \text { p. } 10.
\]

\[
\mathrm{K} (x _ {i}), \mathrm{K} (\mathrm{x}), \mathrm{K} (x _ {1}, \dots , x _ {n}): \mathrm{V}, \text { p. } 10.
\]

\[
h (\mathrm{L}), [ \mathrm{A}: \mathrm{K} ] _ {s}, \mathscr {H} (\mathrm{A}): \mathrm{V}, \text { p. } 31.
\]

\[
\mathrm{E} _ {s}: \mathrm{V}, \text { p. } 44.
\]

\[
[ \mathrm{E}: \mathrm{K} ] _ {s}, [ \mathrm{E}: \mathrm{K} ] _ {i}: \mathrm{V}, \text { p. } 31 \text { and } 46.
\]

\[
\mathrm{N} _ {\mathrm{A} / \mathrm{K}} (x), \operatorname{Tr} _ {\mathrm{A} / \mathrm{K}} (x), \mathrm{D} _ {\mathrm{A} / \mathrm{K}} (x _ {1}, \dots , x _ {n}): \mathrm{V}, \text { p. } 47.
\]

\[
\operatorname{Gal} (\mathrm{N} / \mathrm{K}): \mathrm{V}, \text { p. } 58.
\]

\[
k (\Delta), g (E): V, p. 67.
\]

\[
\mathbf {K} _ {ab}: \mathbf {V}, \text { p. } 77.
\]

\[
\mu_ {m} (\mathbf {K}), \mu_ {\infty} (\mathbf {K}), \mathbf {Z} [ 1 / p ]: \mathbf {V}, \text {   p.   } 78.
\]

\[
\mu_ {l\infty} (\mathbf {K}): \mathbf {V}, \text { p. } 79.
\]

\[
\varphi (n): \mathrm{V}, \text { p. } 79.
\]

\[
\mathrm{R} _ {n} (\mathrm{K}), \Phi_ {n}, \chi_ {n}: \mathrm{V}, \text { p. } 81.
\]

\[
\mathrm{K} \left(\mathrm{A} ^ {1 / n}\right), \langle \sigma , a \rangle : \mathrm{V}, \text { p. } 88.
\]

\[
\wp , \mathrm{K} (\wp^ {- 1} (\mathrm{A})), [ \sigma , a \rangle : \mathrm{V}, \text { p. } 91.
\]

\[
\mathbf {F} _ {q} (\Omega), \mathbf {F} _ {q}: \mathbf {V}, \text { p. } 95.
\]

\[
\mathbf {Z} _ {l}, \hat {\mathbf {Z}}: \mathbf {V}, \mathrm{p.} 96.
\]

\[
\sigma_ {q}: \mathrm{V}, \text { p. } 97.
\]

\[
\varphi_ {n}: \mathrm{V}, \text { p. } 97.
\]

deg.tr \(_{K}\) E：第五卷，第110页。

\[
f ^ {\Delta}: \mathbf {V}, \text { p. } 127.
\]

\(\chi(f), f\) 线性映射：第五卷，第132页。

\[
x \mid y, x \mid y: \mathrm{VI}, \text { p. } 5.
\]

\[
(x): \mathrm{VI}, \text { p. } 6.
\]

\[
x \equiv x ^ {\prime} (\mathrm{mod} y): \mathrm{VI}, \text { p. } 6.
\]

\(\sup_{F}(\mathbf{x}_{i})\) （F 是序集 E 的子集）：VI，第 8 页。

\(\gcd (x_{i}), \operatorname{lcm}(x_{i}) : \mathrm{VI}, \text{p. 8.}\)

\(x^{+}, x^{-}, |x|\) （x 是格序群的元素）：VI，第 12 页。

\[
\operatorname{sgn} (x): \mathrm{VI}, \text { p. } 20.
\]

\(\sqrt{a}\) ( \(a\) 一个元素 \(\geqslant 0\) 有序域的）VI，第24页。

\(|z|\) (z 是 \(K(i)\) ，其中 \(K\) 的一个元素，且 \(i^2 = -1\) ）：VI，第27页。

\(\mathbf{GL}^{+}\) （E）（E为定向向量空间）：VI，第29页。

\(M(\alpha), M_{\alpha}\) (M 是 A 上的一个模， \(\alpha \in A\) ）：第七卷，第7页。

\((\mathbf{Z} / n\mathbf{Z})^{*}, \mathrm{U}(p^{n})\) ：第 VII 卷，第 12 页。

\(c_{\mathsf{L}}(x)\) ，设 L 为主理想整环 A 上的自由模， \(x\in \mathbf{L}\) ：VII，第16页。

\(m(0), m(p^n)\) ( \(p\) 不可约的， \(n\) 为整数，并假设在 \(\geqslant 1\) ）：VII，第24页。

D(M)（M是主理想整环A上的模， \(N^{0}\) ：VII，第25至27页。

\[
c _ {\mathbf {M}}: \mathbf {M} \rightarrow \mathbf {D} (\mathbf {D} (\mathbf {M})): \text { VII,   p. } 26.
\]

\(\mathbf{M}_u\) （M是模， \(u\) 的一个自同态， \(\mathbf{M}\) ）：VII，第28页。

\(V_{\alpha}\) （V是向量空间， \(\alpha \in \mathbf{K}\) ）：VII，第30页。

\(\chi_{u}, \chi_{U}\) ：VII，第30页。

\(u_{s}, u_{n}\) （u是自同态）：VII，第43和45页。

\(f_{s}, f_{u}\) ( \(f\) 一个自同构）：VII，第46页。

\BourbakiUnnumberedChapter{术语索引}{术语索引}

阿贝尔闭包：V，第77页。

阿贝尔扩张：V，第77页。

绝对半单自同态：VII，第42页。

不等式相加：VI，第2页。

伴随（通过……得到的扩张）：V，第11页。

代数闭包：V，第22页。

代数中的代数元素：V，第16页。

代数扩张：V，第17页。

代数数：V，第20页。

代数闭域：V，第20页。

代数相关族、子集：V，第106页。

代数不相交扩张：V，第112页。

代数自由扩张族：V，第115页。

代数自由族、子集：IV，第4页，V，第106页。

阿基米德有序群：VI，第35页，习题31。

阿廷定理：V，第65页。

阿廷-施赖埃尔定理：VI，第22页。

阿廷-施赖埃尔理论：V，第91页。

相伴元素：VI，第5页。

贝祖恒等式：VII，第2页。

布丹-傅里叶法则：VI，第42页，习题21。

域的特征指数：V，第7页。

环的特征：V，第2页。

特征子模：VII，第67页，习题16。

中国剩余定理：VI，第33页，习题25。

闭的（域，代数上）：V，第20页。

闭的（域，可分上）：V，第45页。

闭半直线：VI，第28页。

闭的（相对代数上）域：V，第19页。

（阿贝尔）闭包：V，第77页。

（代数）闭包：V，第22页。

（代数可分）闭包：V，第45页。

（完全）闭包：V，第5页。

（相对代数）闭包：V，第19页。

闭包（相对p-根）：V，第25页。闭包（相对可分代数）：V，第43页。形式幂级数的系数：IV，第24页。多项式的系数：IV，第1页。与群、幺半群结构相容（序关系）：VI，第1页。与环结构相容（序关系）：VI，第19页。与交换幺半群结构相容（预序关系）：VI，第3页。互补定向：VI，第29页。复合扩张：V，第12页。级数的复合：IV，第97页，习题15。Ω中的共轭类：V，第52页。共轭元素：V，第52页；共轭扩张：V，第52页。常数多项式、常数项：IV，第1页。内容：VII，第16页。互素元素：VI，第13页。TS(M)中的余积：IV，第50页。三次扩张：V，第10页。循环扩张：V，第85页。循环向量空间（关于u）：VII，第29页。分圆扩张：V，第81页。分圆多项式：V，第81页。

可分解模：VII，第23页。

一个多项式的（泛）分解代数：IV，第73页。

分解定理：VI，第11页。

戴德金定理：V，第27页。

扩张的（不可分）次数：V，第46页。

域上代数的次数：V，第10页。

代数元素的次数：V，第16页。

多项式的（全）次数：IV，第3页。

有理分式的（全）次数：IV，第20页。

扩张的（可分）次数：V，第31页。

高度为 \(\leqslant 1\) 的p-根扩张中的（偏）导子：V，第103页。

形式幂级数的（偏）导数：IV，第32页。

多项式的（偏）导数：IV，第6页。

笛卡尔符号法则：VI，第42页，习题21。

对角、可对角化的自同态集：VII，第40页。

可对角化、被对角化的代数：V，第28页。

直接基：VI，第29页。

方向向量：VI，第28页。

半直线的方向：VI，第28页。

首一多项式的判别式：IV，第81页。

多项式的判别式：IV，第83页。

元素序列的判别式：V，第47页。

（代数）不相交的扩张：V，第112页。

（线性）不相交的扩张：V，第14页。

TS(M)中的除幂：IV，第45页。

级数的除幂：IV，第96页，习题13。

整除关系：VI，第5页。

（欧几里得）除法：IV，第10页。

（最大公）因子：IV，第12页，VI，第8页。

元素的因子：VI，第5页。

模的（初等）因子：VII，第24页。

二重根：IV，第15页。

特征空间：VII，第30页。

特征值：VII，第30页。

模的初等因子：VII，第24页。

初等对称多项式：IV，第61页。

平展代数：V，第28页。

欧几里得引理：VI，第15页。

欧几里得整环：VII，第49页，习题7。

欧拉恒等式：VI，第8页，IV，第89页，习题6。

欧拉-拉格朗日定理：VI，第26页。

有理分式在原点的展开：IV，第30页，V，第39页。

形式幂级数的指数：IV，第39页。

（指数型）序列：IV，第91页，习题1。

域的扩张：V，第9页。

（代数恒等式的）扩张原理：IV，第18页。

有限扩张：V，第10页。

有限生成扩张：V，第11页。

n次型：IV，第2页。

形式幂级数：IV，第24页。

（广义）形式幂级数：IV，第38页。

（有理）分式域：IV，第19页。

弗罗贝尼乌斯同态：V，第4页，V，第92页。

伽罗瓦扩张：V，第57页。

扩张、多项式的伽罗瓦群：V，第58页。

伽罗瓦理论：V，第67页。

模的伽马代数：IV，第92页，习题2。

高斯整数：VII，第1页。

（生成的）扩张：V，第11页。

（生成的）伽罗瓦扩张：V，第57页。

（生成的）拟伽罗瓦扩张：V，第55页。

扩张的生成族：V，第11页。

元素的最大公因数（GCD）：VI，第8页。

多项式的最大公因数：IV，第12页。

主理想的最大公因数：VII，第3页。

半直线（闭的，开的）：VI，第28页。

半直线（相反的）：VI，第28页。

汉克尔行列式：IV，第90页，习题1。

 \(p\) -根元素：V，第24页。

 \(p\) -根式扩张的高度：V，第26页。

 \(\mathbf{K}(\mathbf{T})\) 中元素的高度：V，第148页，习题11。

埃尔米特插值公式：IV，第88页，习题13。

埃尔米特简化型：VII，第68页，习题21。

形式幂级数的齐次分量：IV，第25页。

不完全域：V，第7页。

不完全度：V，第170页，习题1。

不可分解模：VII，第23页。

未定元：IV，第1页。

线性映射的指标：V，第132页。

欧拉函数：V，第79页。

不可分元素：VII，第17页。

膨胀同态：V，第70页。

不可分次数：V，第46页。

整理想：VI，第6页。

线性映射的不变因子：VII，第22页。

模的不变因子：VII，第20页。

子模的不变因子：VII，第19页。

自同态的不变量（相似性）：VII，第32页。

不可约元素：VI，第17页。

不可约多项式：IV，第13页。

等压多项式：IV，第3页。

若尔当分解：VII，第43页。

若尔当分解（乘性的）：VII，第46页。

若尔当矩阵：VII，第35页。

库默尔扩张：V，第162页，习题7。

库默尔理论：V，第88页。

拉格朗日插值公式：IV，第16页。

运算（多项式的）：IV，第94页，习题9。

首项系数：IV，第3页。

元素的最小公倍数（LCM）：VI，第8页。

主理想的最小公倍数：VII，第3页。

莱夫谢茨原理：V，第117页。

有序群的字典序积：VI，第7页。

线性拓扑：IV，第28页。

线性不相交扩张：V，第13页。

线性拓扑化代数：IV，第28页。

形式幂级数的对数：IV，第40页。对数导数：IV，第41页。吕罗特定理：V，第149页，习题11。

麦克莱恩准则：V，第43页，V，第123页。极大有序域：VI，第25页。元素的极小多项式：V，第16页。自同态的极小多项式：VII，第29页。闵可夫斯基引理：VII，第50页，习题9。首一多项式：IV，第3页。单生成扩张：V，第11页。单项式：IV，第1页。重因子（无重因子的多项式）：IV，第14页。元素的倍数：VI，第5页。重根：IV，第15页。特征值的（几何）重数：VII，第30页。根的重数：IV，第15页。初等因子的重数：VII，第24页。

负元素：VI，第4页。

负的 \(n\) -向量：VI，第29页。

负基：VI，第29页。

牛顿多边形：V，第150页，习题2。

牛顿关系式：IV，第70页，IV，第75页。

幂零分量：VII，第45页。

正规基，正规基定理：V，第73页。

正规扩张：V，第53页。

范数：V，第47页。

开半直线：VI，第28页。

相反半直线：VI，第28页。

（广义）形式幂级数的阶：IV，第25页，IV，第38页。

根的阶：IV，第15页。

可序域：VI，第39页，习题8。

有序扩张：VI，第21页。

有序域：VI，第20页。

有序群、幺半群：VI，第1页。

有序环：VI，第19页。

向量空间的定向：VI，第29页。

定向仿射空间、向量空间：VI，第29页。

\(p\) -进整数（环）：V，第96页。

\(p\) -根闭包：V，第25页。

\(p\) -根元素：V，第24页。

\(p\) -根扩张：IV，第19页。

\(p\) -挠群：VII，第10页。

\(\pi\) -准素分量：VII，第8页。

\(\pi\) -准素模：VII，第7页。

环的完全闭包：V，第5页。

完全域：V，第7页。

特征p的完全环：V，第5页。

多项式，多项式代数：IV，第1页。

多项式函数：IV，第4页。

多项式律：IV，第94页，习题9。

多项式映射：IV，第57页。

多项式映射（齐次）：IV，第55页。

正基：VI，第29页。

正元素：VI，第4页。

正 \(n\) -向量：VI，第29页。

幂级数（广义形式）：IV，第24页，IV，第38页。

预序群、幺半群：VI，第3页。

准素扩张：V，第139页。

素域、子域：V，第1页。

本原元素（定理）：V，第40页。

本原单位根：V，第81页。

主分式理想：VI，第6页。

主理想整环：VII，第1页。

主理想环：VII，第49页，习题6。

有序群的（字典序）积：VI，第7页。

有序群的积：VI，第7页。

积定向：VI，第29页。

Puiseux定理：V，第150页，习题2。

纯基：V，第106页。

纯扩张：V，第106页。

纯子模：VII，第55页，习题7。

毕达哥拉斯域：VI，第39页，习题8。

二次扩张：V，第10页。

二次互反律（定理）：V，第166页，习题23。

拟伽罗瓦扩张：V，第53页。

欧几里得除法中的商：IV，第11页。

商定向：VI，第29页。

有理分式：IV，第19页。

有理函数：IV，第21页。

既约环：V，第34页。

正则代数：V，第140页。

正则扩张：V，第141页。

域在扩张中的相对代数闭包：V，第19页。

在扩张域中相对代数封闭的域：V，第19页。互素多项式：IV，第12页。欧几里得除法中的余数：IV，第11页。不可约元素的（完全）代表系：VII，第3页。限制同态：V，第60页，V，第70页。两个多项式的结式：IV，第76页。多项式的根：IV，第14页。单位根：V，第78页。符号法则：VI，第19页。

（绝对）半单分量：VII，第45页。（绝对）半单自同态：VII，第42页。半单自同态：VII，第41页。半单模：VII，第9页。可分代数：V，第119页。可分代数闭包（相对）：V，第44页。可分代数扩张：V，第36页。可分闭包：V，第44页。可分次数：V，第31页。可分元素：V，第39页。可分扩张：V，第121页。可分多项式：V，第37页。可分闭域：V，第45页。可分超越基：V，第136页。E上取值于F的级数：IV，第96页，习题11。（Γ-）集：V，第75页。元素的符号：VI，第20页。相似自同态、矩阵：VII，第29页。自同态的相似不变量：VII，第32页。单模：VII，第25页。单根：IV，第15页。分裂域、扩张：V，第21页。子扩张：V，第9页。子域：V，第1页。有理分式中可代入的元素：IV，第21页。形式幂级数中的代入：IV，第29页。多项式中的代入：IV，第4页。有理分式中的代入：IV，第21页。（初等）对称多项式：IV，第61页。对称形式幂级数：IV，第67页。对称多项式：IV，第61页。对称张量的对称积：IV，第43页。对称有理分式：IV，第67页。对称张量：IV，第42页。对称化张量：IV，第43页。

切线性映射：IV，第36页。形式幂级数的泰勒公式：IV，第31页。

多项式的泰勒公式：IV，第 8 页。形式幂级数的项、常数项：IV，第 24 页。多项式的项、常数项：IV，第 1 页。迹：V，第 47 页。超越基：V，第 109 页。超越次数：V，第 110 页。超越元：V，第 16 页。超越扩张：V，第 17 页。平移同态：V，第 70 页。可三角化自同态：VII，第 35 页。平凡扩张：V，第 9 页。

乌尔姆–齐平定理：VII，第 59 页，习题14。

幂单分量：VII，第 46 页。

幂单自同态：VII，第 46 页。

环的单位：VI，第 5 页。

首一多项式的通用分解：IV，第 73 页。

权：IV，第 3 页。

威尔逊公式：V，第 94 页。

多项式的零点：IV，第 14 页。
